\providecommand{\toplevelprefix}{../..}
\documentclass[../../book-main_ko.tex]{subfiles}

\begin{document}

\chapter{펼쳐진 최적화로서의 심층 표현}

\section{펼쳐진 최적화를 통한 화이트박스 심층 네트워크}\label{sec:chap4-white-box-model-via-unrolling}

이제, \Cref{sec:chap4-representation-learning-problem}에서 논의한 바와 같이 레이트 감소 또는 정보 이득을 최대화하는 것이 원하는 표현으로 이어진다는 데 동의한다면, 남은 질문은 데이터 $\X$에서 최적의 표현 $\Z^*$로 가는 (비선형) 매핑을 어떻게 구성하고 학습하는가이다. 이는 데이터의 기본 구조를 효과적으로 포착하고 최적의 표현을 충실하게 실현할 수 있는 네트워크 아키텍처와 학습 알고리즘을 설계하는 것을 포함한다.


\subsection{펼쳐진 경사 하강으로부터의 심층 네트워크}

이전 장에서, 우리는 데이터의 선형 판별 표현을 학습하기 위한 원칙적인 목표로 레이트 감소 목표 \eqref{eqn:maximal-rate-reduction}를 제시했다. 그러나 우리는 입력 데이터 $\x$에서 그러한 표현을 추출하기 위한 특징 매핑 $\z = f(\x, \bm \theta)$의 아키텍처를 명시하지 않았다. 
직선적인 선택은 ResNet과 같은 기존의 심층 네트워크를 사용하여 $f(\x, \bm \theta)$를 구현하는 것이다. \Cref{eg:Rate-Reduction-CIFAR10}에서 보았듯이, 이러한 선택은 종종 경험적으로 괜찮은 성능으로 이어진다. 그럼에도 불구하고, 임의의 심층 네트워크를 채택하는 데에는 몇 가지 해결되지 않은 문제가 남아있다. 학습된 특징 표현이 이제 더 해석 가능해졌지만, 네트워크 자체는 여전히 {\em 그렇지 않다}. 선택된 "블랙박스" 네트워크가 원하는 MCR$^2$ 목표를 최적화할 수 있는 이유가 명확하지 않다. 좋은 경험적 결과(예: ResNet 사용)가 네트워크의 아키텍처 및 연산자에 대한 특정 선택을 반드시 정당화하는 것은 아니다: 왜 심층 계층 모델이 필요한가; 추가 계층은 무엇을 개선하거나 단순화하려고 하는가; 얼마나 넓고 깊어야 적절한가; 또는 사용된 합성곱(인기 있는 다중 채널 형태) 및 비선형 연산자(예: ReLU 또는 소프트맥스)에 대한 엄격한 정당성이 있는가? 

이 장에서는 \eqref{eqn:maximal-rate-reduction}에 정의된 레이트 감소 $\Delta R_{\epsilon}(\Z \mid \bm \Pi)$를 최대화하기 위해 경사 상승법을 사용하는 것이 원하는 매핑을 실현하는 "화이트박스" 심층 네트워크로 자연스럽게 이어진다는 것을 보여준다. 모든 네트워크 계층, 선형/비선형 연산자 및 매개변수는 {\em 순전히 순방향 전파 방식으로 명시적으로 구성된다}. 더욱이, 이러한 네트워크 아키텍처는 기존의 경험적으로 설계된 심층 네트워크와 유사하여 그 설계에 대한 원칙적인 정당성을 제공한다.

\paragraph{코딩 레이트 감소를 위한 경사 상승법.} 이전 장에서, 우리는 선형 판별 표현(LDR)을 찾기 위해 수학적으로 데이터 $\X = [\x_1, \ldots, \x_N] \in \Re^{D \times N}$(또는 데이터에서 추출한 초기 특징\footnote{다음 섹션에서 이러한 특징 추출의 필요성을 보게 될 것이다.})에서 다음 코딩 레이트 감소 목표를 최대화하는 최적의 표현 $\Z = [\z_1, \ldots, \z_N] \in \Re^{d \times N}$으로 가는 연속적인 매핑 $f(\cdot): \x \mapsto \z$을 찾고 있음을 알 수 있다:
\begin{equation}\label{eq:mcr-formulation}
\begin{split}
\Delta R_{\epsilon}(\Z \mid \bm{\Pi}) \doteq \underbrace{\frac{1}{2}\log\det \Big(\I + {\alpha} \Z \Z^{\top} \Big)}_{R_{\epsilon}(\Z)} \;-\; \underbrace{\sum_{k=1}^{K}\frac{\gamma_{k}}{2} \log\det\Big(\I + {\alpha_{k}} \Z \bm{\Pi}_{k} \Z^{\top} \Big)}_{R_{\epsilon}^c(\Z \mid \bm \Pi)},
\end{split}
\end{equation}
여기서 $\epsilon > 0$은 미리 지정된 양자화 오차이고 간단함을 위해 다음과 같이 표기한다\footnote{\Cref{ch:compression}에 비해 약간 단순화된 표기법을 사용함에 유의하라.}
\begin{align*}
    \alpha \doteq \frac{d}{N\epsilon^2}, \qquad \alpha_{k} \doteq \frac{d}{\mathrm{tr}(\bm{\Pi}_{k})\epsilon^2}, \qquad \gamma_{k} \doteq \frac{\mathrm{tr}(\bm{\Pi}_{k})}{N}, \qquad \text{for}\ k = 1,\ldots, K.
\end{align*}

질문은 $\bm x$에서 $\bm z$로 가는 그러한 연속적인 매핑 $f(\cdot,\bm \theta)$를 찾는 {\em 구성적인} 방법이 있는지 여부로 귀결된다. 이를 위해, $\Z \subseteq \mathbb{S}^{d-1}$의 함수로서 목표 $\Delta R_{\epsilon}$를 점진적으로 최대화하는 것을 고려해 보자. 선택할 수 있는 많은 최적화 기법이 있을 수 있지만, 간단함을 위해 먼저 가장 간단하다고 할 수 있는 투영 {\em 경사 상승법}(PGA) 기법을 고려한다:\footnote{$\Z_j$의 아래첨자 $j$는 $j$번째 클래스의 특징을 나타내고 $\Z^{\ell}$의 위첨자 $\ell$은 $\ell$번째 반복 또는 계층에서의 모든 특징을 나타냄에 유의하라.} 
\begin{equation}
\bm Z^{\ell+1}   \; \propto \; \bm Z^{\ell} + \eta \cdot \frac{\partial \Delta R_{\epsilon}}{\partial \bm Z}(\Z^{\ell})
\quad \mbox{s.t.} \quad \Z^{\ell+1} \subseteq \mathbb{S}^{d-1}, \quad \ell = 1, 2, \ldots,
\label{eqn:gradient-descent}
\end{equation}
어떤 스텝 크기 $\eta >0$에 대해, 반복은 주어진 데이터 $\bm Z^{0} = \bm X$로 시작한다.\footnote{다시 말하지만, 간단함을 위해 여기서 초기 특징 $\bm Z^{1}$은 데이터 자체라고 먼저 가정한다. 여기서 $\ell$은 반복 횟수를 나타낸다. 따라서 데이터와 특징은 동일한 차원 $d$를 갖는다. 하지만 이것이 반드시 그럴 필요는 없다. 다음 섹션에서 보게 될 것처럼, 초기 특징은 시작하기 위해 데이터의 어떤 (향상된) 특징일 수 있으며 원칙적으로 다른 (훨씬 더 높은) 차원을 가질 수 있다. 모든 후속 반복은 동일한 차원을 갖는다.}
이 기법은 현재 특징 $\Z^\ell$의 위치를 점진적으로 어떻게 조정해야 결과적인 $\Z^{\ell +1}$이 \Cref{fig:gradient-flow}에서 설명한 것처럼 레이트 감소 $\Delta R_{\epsilon}$를 개선하는지로 해석될 수 있다.
\begin{figure}
\centering
    \includegraphics[width=0.85\linewidth]{\toplevelprefix/chapters/chapter4/figs/redu_gradient_diagram.png} 
    \caption{각 클래스의 데이터를 부분 공간으로 평탄화하고 다른 클래스를 밀어내기 위한 경사 흐름을 통한 점진적 변형.} 
    \label{fig:gradient-flow}
\end{figure} 


간단한 계산은 그래디언트 ${\partial \Delta R_{\epsilon}}/{\partial \bm Z}$가 $\Delta R_{\epsilon}$의 두 항의 다음 도함수를 평가해야 함을 보여준다:
\begin{equation}\label{eqn:expand-directions}
    \frac{1}{2}\frac{\partial \log \det (\I \!+\! \alpha \Z \Z^{\top} )}{\partial \bm Z}(\Z^\ell) = \underbrace{\alpha(\I \!+\! \alpha\Z^\ell (\Z^\ell)^\top)^{-1}}_{\bm E^{\ell} \; \in \Re^{d\times d}}\Z^\ell,
\end{equation}

\begin{equation}\label{eqn:compress-directions}
\frac{1}{2}\frac{\partial \left( \gamma_{k}  \log \det (\I + \alpha_{k} \Z \bm \Pi_k \Z^{\top} )  \right)}{\partial \bm Z}(\Z^\ell) = \gamma_{k}  \underbrace{ \alpha_{k}  (\I +  \alpha_{k} \Z^\ell \bm \Pi_k (\Z^\ell)^{\top})^{-1}}_{\bm C^{\ell}_k \; \in \Re^{d\times d}} \Z^{\ell} \bm \Pi_k.
\end{equation}
위에서, 행렬 $\bm E^{\ell}$은 $\Z^{\ell}$에만 의존하며 모든 특징을 {\em 확장}하여 전체 코딩 레이트를 증가시키는 것을 목표로 한다; 행렬 $\bm C^{\ell}_{k}$는 $k$-클래스의 특징에 의존하며 그것들을 {\em 압축}하여 각 클래스의 코딩 레이트를 줄이는 것을 목표로 한다. 
그러면 완전한 그래디언트 $\frac{\partial \Delta R_{\epsilon}}{\partial \bm Z}(\Z^\ell) \in \Re^{d\times N}$은 다음과 같은 형태이다:
\begin{equation}
\frac{\partial \Delta R_{\epsilon}}{\partial \bm Z}(\Z^\ell)  = \underbrace{\bm E^{\ell}}_{\text{확장}} \Z^{\ell} \;-\; \sum_{k=1}^K \gamma_{k} \underbrace{\bm C_{k}^{\ell}}_{\text{압축}}  \Z^{\ell} \bm{\Pi}_k.
\label{eqn:DR-gradient}
\end{equation}

\begin{remark}[선형 연산자로서의 $\bm E^\ell$과 $\bm C_j^\ell$의 해석]\label{rem:regression-interpretation} 
임의의 $\bm z^\ell \in \mathbb{R}^d$에 대해,
\begin{gather}
    \bm E^\ell \bm z^\ell = \alpha(\bm z^\ell - \bm Z^\ell \bm q^\ell_\star), \qquad
    \mbox{where} \qquad \bm q^\ell_\star \doteq \argmin_{\bm q^\ell} \big\{\alpha \|\bm z^\ell - \bm Z^\ell \bm q^\ell\|_2^2 + \|\bm q^\ell\|_2^2\big\}.
\end{gather}
$\bm q^\ell_\star$는 모든 데이터 포인트 $\bm Z^\ell$에 의한 릿지 회귀의 해와 정확히 일치한다는 점에 유의하라. 따라서, $\bm E^\ell$(그리고 유사하게 $\bm C^\ell_k$도)는 대략적으로(즉, $N$이 충분히 클 때) $\bm Z^\ell$의 열에 의해 생성된 부분 공간의 직교 보완 공간으로의 투영이다. 행렬 $\bm E^\ell$을 해석하는 또 다른 방법은 공분산 행렬 $\Z^\ell (\Z^\ell)^\top$의 고유값 분해를 통하는 것이다. $\Z^\ell (\Z^\ell)^\top \doteq \bm U^\ell \bm \Lambda^\ell (\bm U^\ell)^\top$이고 여기서 $\bm \Lambda^\ell \doteq \diag\left(\lambda^\ell_1, \ldots, \lambda^\ell_d \right)$라고 가정하면, 우리는 다음을 얻는다.
\begin{equation}
\bm E^\ell = \alpha \bm U^\ell\, \diag\left(\frac{1}{1+\alpha\lambda^\ell_1}, \ldots, \frac{1}{1+\alpha\lambda^\ell_d}\right) \left(\bm U^\ell\right)^\top.
\end{equation}
따라서, 행렬 $\bm E^\ell$은 큰 분산 방향은 축소되고 소실 분산 방향은 유지되는 방식으로 벡터 $\bm z^\ell$에 작용한다. 이것들이 바로 특징을 이동시켜 전체 볼륨이 확장되고 코딩 레이트가 증가하도록 하는 방향 \eqref{eqn:expand-directions}이며, 따라서 양의 부호를 갖는다. 반대로, \eqref{eqn:compress-directions}와 관련된 방향은 각 클래스의 특징들이 속해야 할 부분 공간에서 벗어난 "잔차"이다. 이것들이 바로 특징들이 각자의 부분 공간으로 다시 압축되어야 하는 방향이며, 따라서 음의 부호를 갖는다(\Cref{fig:regression-interpretation} 참조). 

\begin{figure}[t]
    \centering
    \includegraphics[width=0.65\linewidth]{\toplevelprefix/chapters/chapter4/figs/expand_compress.png}
    \caption{\small $\bm C^\ell_k$와 $\bm E^\ell$의 해석: $\bm C^\ell_k$는 특징을 저차원 부분 공간으로 축소하여 각 클래스를 압축한다; $\bm E^\ell$은 다른 클래스의 특징들을 대조하고 밀어내어 모든 특징을 확장한다.}
    \label{fig:regression-interpretation}
    \vspace{-0.1in}
\end{figure}


본질적으로, 레이트 감소를 위한 경사 상승법의 선형 연산 $\bm E^\ell$과 $\bm C_k^\ell$는 "자동 회귀"를 수행하는 훈련 데이터에 의해 결정된다. 과매개변수화 설정에서의 릿지 회귀에 대한 최근의 새로운 이해 \cite{yang2020rethinking,Wu2020OnTO}는 회귀 변수로 겉보기에 중복적으로 샘플링된 데이터(각 부분 공간에서)를 사용하는 것이 과적합으로 이어지지 않음을 나타낸다.
\end{remark}

\paragraph{경사 유도 특징 맵 증분.} 위에서 경사 상승법은 모든 특징 $\Z^{\ell} = [\z^{\ell}_{1}, \dots, \z^{\ell}_{N}]$을 자유 변수로 간주한다는 점에 유의하라. 증분 $\Z^{\ell+1} - \Z^{\ell} = \eta \frac{\partial \Delta R_{\epsilon}}{\partial \bm Z}(\Z^\ell)$은 아직 전체 특징 도메인 $\z^\ell \in \Re^d$에 대한 변환을 제공하지 않는다. \eqref{eqn:DR-gradient}에 따르면, 그래디언트는 소속이 알려지지 않은 지점에서 평가될 수 없다. \Cref{fig:gradient-flow}에서 설명한 것처럼. 따라서 최적의 $f(\x,\bm  \theta)$를 명시적으로 찾기 위해, 우리는 $\ell$번째 계층 특징 $\z^\ell$에 대한 작은 증분 변환 $g(\cdot, \bm{\theta}^{\ell})$을 구성하여 위의 (투영된) 경사 기법을 모방하는 것을 고려할 수 있다:
\begin{equation}
\z^{\ell + 1}   \; \propto \; \z^{\ell} + \eta\cdot  g(\z^{\ell}, \bm{\theta}^{\ell}) \quad \mbox{subject to} \quad \z^{\ell +1} \in \mathbb{S}^{d-1}
\label{eqn:gradient-descent-transform}
\end{equation}
$\big[g(\z_1^{\ell}, \bm \theta^{\ell}), \ldots, g(\z_N^{\ell}, \bm \theta^{\ell}) \big] \approx \frac{\partial \Delta R_{\epsilon}}{\partial \bm Z}(\Z^\ell)$가 되도록. 즉, 우리는 모든 (훈련) 특징 $\{\z_i^\ell\}_{i=1}^N$을 국소적으로 변형하는 경사 흐름 $\frac{\partial \Delta R_{\epsilon}}{\partial \bm Z}$을 전체 특징 공간 $\z^\ell \in \Re^d$에 정의된 연속 매핑 $g(\z,\bm \theta)$로 근사해야 한다.  
증분 \eqref{eqn:gradient-descent-transform}을 연속 미분 방정식의 이산화된 버전으로 해석할 수 있다는 점에 유의하라:
\begin{equation}
\dot{\z} = g(\z, \theta).
\end{equation}
따라서 이렇게 구성된 (심층) 네트워크는 어떤 신경 ODE \cite{chen2018neural}로 해석될 수 있다. 그럼에도 불구하고, 흐름 $g$가 어떤 일반적인 구조로 선택되는 신경 ODE와 달리, 여기서 우리의 $g(\z, \bm\theta)$는 특징 집합에 대한 레이트 감소의 경사 흐름을 모방하는 것이다(\Cref{fig:gradient-flow}에 표시된 대로): 
\begin{equation*}
    \dot{\Z} = \frac{\partial \Delta R_{\epsilon}}{\partial \bm Z},
\end{equation*} 
그리고 그 구조는 다른 사전 지식이나 휴리스틱 없이 이 목표로부터 전적으로 파생되고 완전히 결정된다.  

그래디언트 \eqref{eqn:DR-gradient}의 구조를 검토하면, 증분 변환 $g(\z^\ell, \bm \theta^\ell)$에 대한 자연스러운 후보가 다음과 같은 형태임을 시사한다:
\begin{equation}
    g(\z^\ell, \bm \theta^\ell) \; \doteq \; \bm E^\ell \z^\ell - \sum_{k=1}^K \gamma_{k} \pi_k(\z^\ell)\bm C_{k}^{\ell}  \z^{\ell}  \in \Re^d,
    \label{eqn:DR-gradient-transform}
\end{equation}
여기서 $\pi_k(\z^\ell) \in [0,1]$은 $\z^{\ell}$이 $k$번째 클래스에 속할 확률을 나타낸다. 증분 맵 매개변수 $\bm \theta^\ell$은 다음에 의존한다: 첫째, 훈련 $\Z^\ell$의 특징 통계에만 의존하는 $\bm E^{\ell}$과 $\{ \bm C^{\ell}_{k}\}_{k=1}^{K}$로 표현되는 선형 맵 집합; 둘째, 임의의 특징 $\z^\ell$의 소속 $\{ \pi_k(\z^\ell)\}_{k=1}^K$.
훈련 샘플 $\Z^\ell$에 대해, 소속 $\bm \Pi_k$가 알려져 있으므로,  이렇게 정의된 $g(\z^\ell, \bm \theta)$는 그래디언트 $\frac{\partial \Delta R_{\epsilon}}{\partial \bm Z}(\Z^\ell)$에 대한 값을 정확히 제공한다는 점에 유의하라.  


우리는 훈련 샘플에 대해서만 소속 정보를 가지고 있으므로, \eqref{eqn:DR-gradient-transform}에 정의된 함수 $g(\cdot)$는 훈련 시에만 평가할 수 있다. $g(\cdot)$를 전체 특징 공간으로 외삽하려면, 두 번째 항에서 $\pi_k(\z^\ell)$를 추정해야 한다. 기존 딥러닝에서는 이 맵이 일반적으로 심층 네트워크로 모델링되고 훈련 데이터로부터 학습된다(예: {\em 역전파}를 통해). 그럼에도 불구하고, 여기서 우리의 목표는 정확한 분류기 $\pi_{k}(\z^\ell)$를 이미 학습하는 것이 아니다. 대신, 우리는 $g(\cdot)$가 그래디언트 $\frac{\partial \Delta R_{\epsilon}}{\partial \bm Z}$를 잘 근사하기 위해 충분히 좋은 클래스 정보 추정치만 필요하다. 



\Cref{rem:regression-interpretation}에서 주어진 선형 맵 $\bm E^\ell$과 $\bm C_k^\ell$의 기하학적 해석으로부터, 항 $\bm p_{k}^{\ell} \doteq \bm C^{\ell}_k \z^{\ell}$은 (대략적으로) $\z^{\ell}$을 각 클래스 $j$의 직교 보완 공간으로 투영한 것으로 볼 수 있다. 따라서, $\|\bm p_{j}^{\ell}\|_2$는 $\z^\ell$이 클래스 $j$에 있으면 작고 그렇지 않으면 크다. 이는 다음 소프트맥스 함수를 기반으로 소속을 추정하도록 동기를 부여한다:
\begin{equation}
    \widehat{\vpi}(\z^\ell) \doteq \softmax\rp{-\lambda\mat{\norm{\vC^{\ell}_{1}\vz^{\ell}}_{2} \\ \vdots \\ \norm{\vC^{\ell}_{K}\vz^{\ell}}_{2}}} = \frac{1}{\sum_{k = 1}^{K}\exp(-\lambda \norm{\vC^{\ell}_{k}\vz^{\ell}}_{2})}\mat{\exp(-\lambda\norm{\vC^{\ell}_{1}\vz^{\ell}}_{2}) \\ \vdots \\ \exp(-\lambda\norm{\vC^{\ell}_{K}\vz^{\ell}}_{2})} \in [0,1]^{K}.
\end{equation}
따라서, \eqref{eqn:DR-gradient-transform}의 두 번째 항은 이 추정된 소속으로 근사될 수 있다:
\begin{align}
\sum_{k=1}^K \gamma_{k} \pi_k(\z^\ell)\bm C_{k}^{\ell}  \z^{\ell} 
\; \approx \;  \sum_{k=1}^K \gamma_{k} \widehat{\pi}_k(\z^\ell)  \bm{C}^{\ell}_k  \z^{\ell} 
\; \doteq \; \bm \sigma\Big([\bm{C}^{\ell}_{1} \z^{\ell}, \dots, \bm{C}^{\ell}_{K} \z^{\ell}]\Big),
\label{eqn:soft-residual}
\end{align}
이는 특징 $\z^\ell$의 출력을 $K$ 그룹의 필터: $[\bm{C}^{\ell}_{1}, \dots, \bm{C}^{\ell}_{K}]$를 통해 비선형 연산자 $\bm \sigma(\cdot)$로 나타낸다. 비선형성은 해당 필터로부터의 특징 응답에 기반한 클래스 소속의 "소프트" 할당으로 인해 발생한다는 점에 유의하라.  

전반적으로, \eqref{eqn:gradient-descent-transform},  \eqref{eqn:DR-gradient-transform}, \eqref{eqn:soft-residual}을 결합하면, 
$\z^{\ell}$에서 $\z^{\ell+1}$로의 증분 특징 변환은 이제 다음과 같이 된다
\begin{equation}\label{eqn:layer-approximate}
\begin{aligned}
\z^{\ell+1}  &\propto \; \z^\ell +  \eta \cdot  \bm E^{\ell} \z^{\ell} - \eta\cdot  \bm \sigma\big([\bm{C}^{\ell}_{1} \z^{\ell}, \dots, \bm{C}^{\ell}_{K} \z^{\ell}]\big)  \\
&= \; \z^\ell +  \eta \cdot g(\z^\ell, \bm \theta^\ell) \qquad \mbox{s.t.} \quad \z^{\ell +1} \in \mathbb{S}^{d-1},
\end{aligned}
\end{equation}
위에서 정의된 비선형 함수 $\bm \sigma(\cdot)$와 모든 계층별 매개변수를 수집하는 $\bm \theta^\ell$과 함께. 즉  $\bm \theta^\ell =\left\{\bm E^\ell, \bm{C}^{\ell}_{1}, \dots, \bm{C}^{\ell}_{K}, \gamma_{k}, \lambda\right\}$이다. 각 계층의 특징은 항상 단위 구 $\mathbb S^{d-1}$에 투영하여 "정규화"되며, $\mathcal P_{\mathbb S^{d-1}}$로 표기된다는 점에 유의하라. \eqref{eqn:layer-approximate}의 증분 형태는 \Cref{fig:arch-ReduNet}(a)의 다이어그램으로 설명될 수 있다.

\begin{figure}[t]
    \begin{subfigure}[t]{0.35\textwidth}
        \centering
        \includegraphics[width=\textwidth]{\toplevelprefix/chapters/chapter4/figs/redunet_layer.png}
        \caption{\textbf{ReduNet}}
    \end{subfigure}
    \hfill 
    \begin{subfigure}[t]{0.6\textwidth}
        \centering
        \includegraphics[width=\textwidth]{\toplevelprefix/chapters/chapter4/figs/resnet_resnext.pdf}
        \caption{\textbf{ResNet} and \textbf{ResNeXt}.}
    \end{subfigure}
    \caption{\small ReduNet의 네트워크 아키텍처 및 다른 아키텍처와의 비교. \textbf{(a)}: 레이트 감소 최적화를 위한 경사 상승법의 한 반복에서 파생된 \textbf{ReduNet}의 계층 구조. \textbf{(b)} (왼쪽): ResNet의 한 계층~\cite{he2016deep}; 및 \textbf{(b)} (오른쪽): ResNeXt의 한 계층~\cite{ResNEXT}. \Cref{sec:shift-invariant}에서 보게 될 것처럼, ReduNet의 선형 연산자 $\bm E^\ell$와 $\bm{C}_k^\ell$는 이동 불변성이 부과될 때 자연스럽게 (다중 채널) 합성곱이 된다.}
    \label{fig:arch-ReduNet}
\end{figure}

\begin{algorithm}[t]
    \caption{ReduNet 훈련 알고리즘}
    \label{alg:redunet_training}
    \begin{algorithmic}[1]
        \Require{$\vX = [\vx_1, \ldots, \vx_N]\in \Re^{D \times N}$, $\bm{\Pi} = \{\bm \Pi_k\}_{k=1}^K$, $\epsilon > 0$, $\lambda$, 및 학습률 $\eta$.}
        \Ensure{학습된 매개변수 \(\{\vE^{\ell}\}_{\ell = 1}^{L}, \{\{\vC^{\ell}_{k}\}_{k = 1}^{K}\}_{\ell = 1}^{L}, \{\gamma_{k}\}_{k = 1}^{k}\)}

        \Procedure{ReduNetTraining}{$\vX, \vPi, \epsilon, \lambda, \eta$}
            \Statex{\quad\ \texttt{\# 상수 정의}}

            \State{\(\alpha \gets d/(N\epsilon^{2})\)}
            \For{\(k \in \{1, \dots, K\}\)}
            \State{\(\alpha_{k} \gets D/(\tr(\vPi_{k})\epsilon^{2})\)}
            \State{\(\gamma_{k} \gets \tr(\vPi_{k})/D\)}
            \EndFor
            \Statex{}

            \Statex{\quad\ \texttt{\# ReduNet 계층별 반복}}
            \State{\(\vZ^{1} = \mat{\vz_{1}^1, \dots, \vz_{N}^1} \gets \vX\)} \Comment{ReduNet 계층별 반복 초기화}
            \For{\(\ell \in \{1, \dots, L\}\)}
                \Statex{\quad\ \texttt{\# 단계 1: 네트워크 매개변수 계산} \(\vE^{\ell}, \{\vC^{\ell}_{k}\}_{k = 1}^{K}\)}
                \State{\(\vE^{\ell} \gets \alpha\left(\vI + \alpha \vZ^{\ell}(\vZ^{\ell})^{\top}\right)^{-1} \in \R^{d \times d}\)}
                \For{\(k \in \{1, \dots, K\}\)}
                    \State{\(\vC^{\ell}_{k} \gets \alpha_{k}\left(\vI + \alpha_{k}\vZ^{\ell}\vPi_{k}(\vZ^{\ell})^{\top}\right)^{-1} \in \R^{d \times d}\)}
                \EndFor
                \Statex{}
                \Statex{\qquad\quad\texttt{\# 단계 2: 특징 업데이트 \(\vZ^{\ell}\)}}
                \For{\(i \in \{1, \dots, N\}\)}
                    \State{\(\hat{\vpi}(\vz^{\ell}_i) \gets \displaystyle \softmax(-\lambda[\norm{\vC^{\ell}_{1}\vz^{\ell}_i}_{2}, \dots, \norm{\vC^{\ell}_{K}\vz^{\ell}_i}_{2}]) \in [0, 1]^{K}\)} \Comment{소프트 할당 계산 \(\hat{\vpi}(\vz^{\ell}_i)\)}
                    \State{\(\displaystyle\vz^{\ell + 1}_i \gets \cP_{\bS^{d - 1}}\rp{\vz^{\ell}_i + \eta \bp{\vE^{\ell}\vz^{\ell}_i - \sum_{k = 1}^{K}\gamma_{k}\hat{\pi}_{k}(\vz^{\ell}_i) \vC^{\ell}_{k}\vz^{\ell}_i}} \in \R^{d}\)} \Comment{\(\vz^{\ell}_i\)로부터 특징 \(\vz^{\ell + 1}_i\) 업데이트}
                \EndFor
            \EndFor
            \State{\Return{\(\{\vE^{\ell}\}_{\ell = 1}^{L}, \{\{\vC^{\ell}_{k}\}_{k = 1}^{K}\}_{\ell = 1}^{L}, \{\gamma_{k}\}_{k = 1}^{K}\)}} \Comment{모든 네트워크 매개변수 반환.}
        \EndProcedure
    \end{algorithmic}
\end{algorithm}


\paragraph{레이트 감소 최적화를 위한 심층 네트워크.} 증분은 레이트 감소 $\Delta R_\epsilon$에 대한 경사 상승법을 모방하도록 구성된다는 점에 유의하라. 따라서 위 과정을 통해 특징을 반복적으로 변환함으로써, 실험 섹션에서 보게 될 것처럼 레이트 감소가 증가할 것으로 기대한다. 이 반복적인 과정은, 예를 들어 $L$번의 반복 후에 수렴하면, 입력 $\x = \z^0$에 대해 원하는 특징 맵 $f(\x, \bm \theta)$를 제공하며, 이는 정확히 {\em 심층 네트워크}의 형태를 띤다. 여기서 각 계층은 \Cref{fig:arch-ReduNet} 왼쪽에 표시된 구조를 갖는다:
\begin{equation}\label{eqn:ReduNet}
\begin{aligned}
f(\x, \bm \theta)\; =&  \;\;f^L \circ f^{L-1} \circ  \cdots \circ f^1 \circ
    f^0(\z^0),  \\ 
f^\ell(\z^\ell, \bm \theta^\ell) \; \doteq & \;\; \z^{\ell+1} = \mathcal{P}_{\mathbb{S}^{n-1}}[\z^{\ell} + \eta\cdot g(\z^{\ell}, \bm \theta^{\ell})], \\
g(\z^{\ell}, \bm \theta^{\ell}) \; =&\;\; \bm E^{\ell} \z^{\ell} -  \bm \sigma\big([\bm{C}^{\ell}_{1} \z^{\ell}, \dots, \bm{C}^{\ell}_{K} \z^{\ell}]\big).
\end{aligned}
\end{equation}
이 심층 네트워크는 레이트 \textbf{감소}(\textbf{redu}ction)를 최대화하는 데서 파생되었으므로, 우리는 그것을 \textbf{ReduNet}이라고 부른다. ReduNet의 아키텍처를 경험적으로 설계된 인기 있는 네트워크인 \textbf{ResNet}과 \textbf{ResNeXt}(\Cref{fig:arch-ReduNet}에 표시됨)와 비교하면, 그 유사성은 다소 기묘하다. 개념적으로, ReduNet은 인기 있는 전문가 혼합(\textbf{MoE}) 아키텍처 \cite{MoE}를 정당화하는 데 사용될 수도 있다. 각 병렬 채널 $\bm{C}^{\ell}_k$는 각 객체 클래스에 대해 훈련된 "전문가"로 볼 수 있기 때문이다.

\begin{figure}[t]
    \centering
    \includegraphics[width=0.45\linewidth]{\toplevelprefix/chapters/chapter4/figs/MoE.png} \hspace{5mm}
    \includegraphics[width=0.45\linewidth]{\toplevelprefix/chapters/chapter4/figs/switched-transformer.png}
    \caption{왼쪽: 전문가 혼합(MoE) 심층 네트워크 \cite{MoE}. 오른쪽: 1.7조 개의 매개변수로 MoE를 구현하는 데 사용되는 희소성 촉진 스위치 트랜스포머 \cite{Fedus-2022}.}
    \label{fig:enter-label}
\end{figure}

우리는 ReduNet의 훈련과 평가를 각각 \Cref{alg:redunet_training}과 \Cref{alg:redunet_evaluation}에 요약한다.  네트워크의 모든 매개변수는 {\em 순방향 전파} 방식으로 계층별로 명시적으로 구성된다는 점에 유의하라. 이 구성에는 역전파가 필요하지 않다! 이렇게 학습된 특징은 분류에 직접 사용될 수 있다(예: 최근접 부분 공간 분류기를 통해).



\begin{algorithm}[!htbp]
    \caption{ReduNet 평가 알고리즘}
    \label{alg:redunet_evaluation}
    \begin{algorithmic}[1]
        \Require{입력 \(\vx \in \R^{D}\), 네트워크 매개변수 \(\{\vE^{\ell}\}_{\ell = 1}^{L}, \{\{\vC^{\ell}_{k}\}_{k = 1}^{K}\}_{\ell = 1}^{L}, \{\gamma_{k}\}_{k = 1}^{K}\), 학습률 \(\lambda\)}
        \Ensure{특징 \(\vz^{L + 1}\)}

        \Procedure{ReduNetEvaluation}{$\vx$}
        \State{\(\vz^{1} \gets \vx \in \R^{D}\)} \Comment{ReduNet 계층별 반복 초기화}
        \For{\(\ell \in \{1, \dots, L\}\)}
        \State{\(\hat{\vpi}(\vz^{\ell}) \gets  \softmax(-\lambda\mat{\norm{\vC^{\ell}_{1}\vz^{\ell}}_{2}, \dots, \norm{\vC^{\ell}_{K}\vz^{\ell}}_{2}}) \in [0, 1]^{K}\)} \Comment{소프트 할당 \(\hat{\vpi}(\vz^{\ell})\) 계산}
        \State{\(\vz^{\ell + 1} \gets \cP_{\bS^{d - 1}}\rp{\vz^{\ell} + \eta \bp{\vE^{\ell}\vz^{\ell} - \sum_{k = 1}^{K}\gamma_{k}\hat{\pi}_{k}(\vz^{\ell}) \vC^{\ell}_{k}\vz^{\ell}}} \in \R^{d}\)} \Comment{\(\vz^{\ell}\)로부터 특징 \(\vz^{\ell + 1}\) 업데이트}
        \EndFor
        \State{\Return{\(\vz^{L + 1}\)}} \Comment{출력 특징 반환}
        \EndProcedure
    \end{algorithmic}
\end{algorithm}


\begin{example}
ReduNet이 특징을 어떻게 변환하는지에 대한 직관을 제공하기 위해, 혼합된 3D 가우시안으로 간단한 예를 제공하고 \Cref{fig:redu-3d-gaussian-diagram}에서 특징이 어떻게 변환되는지 시각화한다. 
$\mathbb{S}^2$에 투영된 $\R^{3}$의 세 가우시안 분포의 혼합을 고려한다. 먼저 3개 클래스에 대한 데이터 포인트를 생성한다: $k=1,2,3$에 대해, $\bm{X}_{k}=[\bm{x}_{k,1}, \ldots, \bm{x}_{k,m}] \in \R^{3\times m}$, $\bm{x}_{k,i} \sim \mathcal{N}(\bm{\mu}_{k}, \sigma_{k}^{2} \I)$, 그리고 ${\pi}(\x_{k,i}) = k$.
$m=500, \sigma_{1}=\sigma_{2}=\sigma_{3}=0.1$로 설정하고, $\bm{\mu}_{1}, \bm{\mu}_{2}, \bm{\mu}_{3} \in \mathbb{S}^2$이다. 
그런 다음 모든 데이터 포인트를 $\mathbb{S}^{2}$에 투영한다. 즉, $\bm{x}_{k,i}/\|\bm{x}_{k,i}\|_{2}$이다. 
네트워크를 구성하기 위해($\ell$-번째 계층에 대한 $\bm{E}^{\ell}, \bm{C}^{\ell}_{k}$ 계산), 반복/계층 수 $L=2,000$, 스텝 크기 $\eta=0.5$, 정밀도 $\epsilon=0.1$로 설정한다. 
우리는 우리의 프레임워크가 수천 개의 계층으로도 안정적인 심층 네트워크로 이어진다는 것을 보여주기 위해서만 이 작업을 수행한다. 
실제로 수천 개의 계층이 필요하지 않을 수 있으며, 새 계층을 추가해도 수익이 감소할 때마다 중단할 수 있다. 
이 예에서는 수백 개의 계층으로 충분하다. 따라서 명확한 최적화 목표는 필요한 네트워크 깊이에 대한 자연스러운 기준을 제공한다.

\Cref{fig:redu-3d-gaussian-diagram}에 표시된 것처럼, 매핑 $f(\cdot, \bm{\theta})$ 후에 동일한 클래스의 샘플이 고도로 압축되어 단일 클러스터로 수렴하고 두 다른 클러스터 사이의 각도가 약 $\pi/2$임을 관찰할 수 있다. 이는 $\mathbb{S}^2$에서 MCR$^2$ 손실의 최적 해 $\Z^{\star}$와 잘 일치한다. 
다른 계층의 특징에 대한 MCR$^2$ 손실은 \Cref{fig:redu-3d-gaussian-diagram}(\textbf{c})에서 찾을 수 있다. 경험적으로, 우리는 구성된 ReduNet이 MCR$^2$ 손실을 최대화할 수 있고 안정적으로 수렴하며 동일한 클래스의 샘플이 하나의 클러스터로 수렴하고 다른 클러스터가 서로 직교한다는 것을 발견했다. 
더욱이, 동일한 분포에서 새 데이터 포인트를 샘플링할 때, 동일한 클래스의 새 샘플이 훈련 샘플과 동일한 클러스터 중심으로 일관되게 수렴하는 것을 발견했다. 
\begin{figure}[t]
    \begin{subfigure}[t]{0.32\textwidth}
        \centering 
        \includegraphics[width=\textwidth]{\toplevelprefix/chapters/chapter4/figs/scatter3d-X_train.pdf}\vspace{-0.1in}
        \caption{$\bm{X}_{\text{train}}$}
    \end{subfigure}
    \hfill
    \begin{subfigure}[t]{0.32\textwidth}
        \centering 
        \includegraphics[width=\textwidth]{\toplevelprefix/chapters/chapter4/figs/scatter3d-Z_train.pdf}\vspace{-0.1in}
        \caption{$\bm{Z}_{\text{train}}$}
    \end{subfigure}
    \hfill
    \begin{subfigure}[t]{0.32\textwidth}
        \centering 
        \includegraphics[width=\textwidth]{\toplevelprefix/chapters/chapter4/figs/scatter3d-loss-traintest.pdf}\vspace{-0.1in}
        \caption{손실 (train/val)}
    \end{subfigure}
    \vspace{-0.1in}
    \caption{\small 3D 가우시안 혼합에 대한 원본 샘플 및 학습된 표현. (\textbf{a})에서 매핑 전 데이터 포인트 $\bm{X}$와 (\textbf{b})에서 매핑 후 학습된 특징 $\bm{Z}$를 산점도로 시각화한다. 각 산점도에서 각 색상은 샘플의 한 클래스를 나타낸다. (\textbf{c})에서는 목적 함수의 값 진행에 대한 플롯도 보여준다.}
    \label{fig:redu-3d-gaussian-diagram}
\end{figure}

\end{example}



\subsection{불변 레이트 감소로부터의 합성곱 네트워크}\label{sec:shift-invariant}

이전 섹션에서는 레이트 감소 목표를 위한 펼쳐진 최적화를 사용하여 심층 네트워크인 ReduNet의 계층별 아키텍처를 유도했다. 
구체적으로, \eqref{eq:mcr-formulation}의 압축 항 $R^c_\epsilon(\Z \mid\bm \Pi)$은 동일한 클래스의 표현을 압축하도록 설계되었다. 그러나 이 공식은 입력 데이터의 가능한 도메인 변환이나 변형을 고려하지 않는다. 예를 들어, 객체를 오른쪽으로 약간 이동해도 이미지의 의미론적 레이블은 변경되지 않는다. 이 섹션에서는 이미지 회전 및 평행 이동과 같은 특정 도메인 변형에 불변하는 레이트 감소 목표를 최대화하여 합성곱 계층을 유도하는 방법을 보여줄 것이다.

많은 클러스터링 또는 분류 작업(예: 이미지에서의 객체 탐지)에서, 우리는 두 샘플이 특정 클래스의 도메인 변형 또는 증강 $\cT = \{\tau \}$에 의해 다른 경우 {\em 동등}하다고 간주한다. 따라서, 우리는 이러한 변형에 {\em 불변}하는 저차원 구조에만 관심이 있다~(즉, 모든 $\tau \in \cT$에 대해 $\x \in \mathcal{M}$인 경우에만 $\tau(\x) \in \mathcal{M}$)\,),
이는 정교한 기하학적 및 위상적 구조를 갖는 것으로 알려져 있으며, 엄격하게 설계된 CNN \cite{Cohen-ICML-2016}으로도 실제 환경에서 정확하게 학습하기 어려울 수 있다.
이 프레임워크에서는 모든 등변 인스턴스가 동일한 부분 공간에 포함되도록 하여 부분 공간 자체가 고려 중인 변환에 불변하도록 하는 매우 자연스러운 방식으로 이를 공식화할 수 있다.

많은 응용 분야, 예를 들어 시리얼 데이터나 영상 데이터에서 데이터의 의미적 의미(레이블)는 특정 변환 $\mathfrak{g} \in \mathbb{G}$(어떤 그룹 $\mathbb{G}$에 대해)에 대해 {\em 불변}이다 \cite{CohenW16,deep-sets-NIPS2017}. 예를 들어, 오디오 신호의 의미는 시간 이동에 불변하며, 이미지 속 객체의 정체성은 이미지 평면에서의 평행 이동에 불변하다. 따라서, 우리는 특징 매핑 $f(\x,\bm \theta)$이 그러한 변환에 대해 엄격하게 불변하기를 선호한다:
\begin{equation}
\mbox{\em 그룹 불변성:}\;   f(\x\circ \mathfrak{g}, \bm \theta) \sim f(\x,\bm \theta),\ \forall \mathfrak{g} \in \mathbb{G},
\end{equation}
여기서 "$\sim$"는 동일한 동치 클래스에 속하는 두 특징을 나타낸다. 불변성이나 등변성을 보장하기 위해 합성곱 연산자가 심층 네트워크에서 일반적인 관행이었지만 \cite{CohenW16}, 평행 이동 및 회전과 같은 간단한 변환에 대해서도 불변성을 {\em 보장}할 수 있는 (경험적으로 설계된) 합성곱 네트워크를 처음부터 훈련시키는 것은 실제로는 여전히 어려운 과제이다 \cite{azulay2018deep,engstrom2017rotation}. 대안적인 접근 방식은 각 계층의 합성곱 필터를 신중하게 설계하여 광범위한 신호에 대한 평행 이동 불변성을 보장하는 것이다. 예를 들어, ScatteringNet \cite{scattering-net} 및 후속 연구 \cite{Wiatowski-2018}에서 웨이블릿을 사용하는 것처럼 말이다. 그러나 일반적인 신호에 대한 불변성을 보장하기 위해 필요한 합성곱의 수는 일반적으로 네트워크 깊이에 따라 기하급수적으로 증가한다. 이것이 바로 이 유형의 네트워크를 보통 몇 개의 계층만으로 깊게 구성할 수 없는 이유이다.

이제, 우리는 MCR$^2$ 원리가 자연스럽고 정확한 방식으로 불변성과 양립할 수 있음을 보여준다: 우리는 모든 변환된 버전 $\{\x\circ \mathfrak{g} \mid \mathfrak{g} \in \mathbb G\}$을 데이터 $\x$와 동일한 클래스로 할당하고 그들의 특징 $\z$를 모두 동일한 부분 공간 $\mathcal S$에 매핑하기만 하면 된다. 따라서 모든 그룹 등변 정보는 부분 공간 내부에만 인코딩되며, 결과적인 부분 공간 집합에 정의된 모든 분류기는 이러한 그룹 변환에 자동으로 불변하게 된다. 1D 회전 및 2D 평행 이동의 예시에 대한 그림은 \Cref{fig:ortho-invariance-diagram}을 참조하라. 다음으로, 그룹 $\mathbb G$가 원형 1D 이동일 때 결과적인 심층 네트워크가 자연스럽게 {\em 다중 채널 합성곱 네트워크}가 됨을 엄격하게 보여줄 것이다.
이렇게 구성된 네트워크는 주어진 데이터 $\X$ 또는 그 특징 $\Z$에 대한 불변성만 보장하면 되기 때문에, 필요한 합성곱의 수는 ScatteringNet과 달리 매우 깊은 네트워크를 통해서도 실제로 일정하게 유지된다.

\begin{figure}[t]
    \begin{subfigure}[t]{0.4\textwidth}
        \centering 
        \includegraphics[width=\textwidth]{\toplevelprefix/chapters/chapter4/figs/ortho_diagram_1d.pdf}
    \end{subfigure}
    \hfill 
    \begin{subfigure}[t]{0.4\textwidth}
        \centering 
        \includegraphics[width=\textwidth]{\toplevelprefix/chapters/chapter4/figs/ortho_diagram_2d.pdf}
    \end{subfigure}
    \caption{이미지 회전(왼쪽) 또는 평행 이동(오른쪽)에 등변/불변하는 탐색된 표현의 그림: 각 클래스의 모든 변환된 이미지는 다른 부분 공간과 비일관적인 동일한 부분 공간으로 매핑된다. 각 부분 공간에 포함된 특징은 변환 그룹에 등변적인 반면 각 부분 공간은 이러한 변환에 불변하다.}\label{fig:ortho-invariance-diagram} 
\end{figure}

\paragraph{1차원 직렬 데이터 및 이동 불변성} 이동에 불변하는 1차원 데이터 $\bm x = [x(0), x(1), \ldots, x(D-1)] \in \Re^D$를 분류하기 위해, 우리는 $\mathbb{G}$를 모든 순환 이동의 그룹으로 삼는다. 각 관측치 $\bm x_i$는 이동된 사본들의 집합 $\{ \x_i \circ \mathfrak{g} \, | \, \mathfrak{g} \in \mathbb G \}$을 생성하며, 이는 다음과 같이 주어진 순환 행렬 $\circm(\bm x_i) \in \Re^{D \times D}$의 열들이다.
\begin{equation}
\circm(\x) \,\doteq\, \left[ \begin{array}{ccccc} x(0) & x(D-1) & \dots & x(2) & x(1) \\ x(1) & x(0) & x(D-1) & \cdots & x(2) \\ \vdots & x(1) & x(0) &\ddots & \vdots \\ x(D-2) &  \vdots & \ddots & \ddots & x(D-1) \\ x(D-1) & x(D-2) & \dots & x(1) & x(0)   \end{array} \right]  \in \Re^{D \times D}.
\end{equation}
순환 행렬의 속성에 대해서는 독자들에게 \cite{Kra2012OnCM}를 참조하도록 한다. 간단히 하기 위해, $\bm Z^1 \doteq [ \z_{1}^1, \dots, \z_{N}^1 ] = \X \in \Re^{d \times N}$이라고 하자.\footnote{다시, 논의를 단순화하기 위해, 초기 특징 $\Z^1$이 $\X$ 자체이므로 동일한 차원 $d$를 갖는다고 가정한다(즉, $D=d$). 그러나 우리가 곧 $\X$를 더 높은 차원으로 끌어올릴 필요가 있음을 보게 될 것이므로 이것이 반드시 그럴 필요는 없다.} 그렇다면 그들의 순환 집합 $\circm(\bm Z^1) = \left[ \circm(\z_{1}^1), \dots, \circm(\z_{N}^1) \right] \in \Re^{d \times (dN)}$으로부터 ReduNet을 구성하면 어떻게 될까? 즉, 우리는 ReduNet에 의해 이 모든 것을 동일한 부분 공간으로 압축하고 매핑하기를 원한다.

이제 데이터 공분산 행렬은 다음과 같다:
\begin{eqnarray}
\circm(\bm Z^1) \circm(\bm Z^1)^\top 
&=& \left[ \circm(\z_{1}^1), \dots, \circm(\z_{N}^1) \right] \left[ \circm(\z_{1}^1), \dots, \circm (\z_{N}^1) \right]^\top \nonumber \\
&=& \sum_{i =1}^N \circm(\z_{i}^1) \circm(\z_{i}^1)^\top \;\in \Re^{d\times d} 
\end{eqnarray}
이 샘플 집합과 관련된 것은 {\em 자동으로} (대칭) 순환 행렬이다. 더욱이, 순환 속성은 합, 역행렬, 곱셈 하에서 보존되기 때문에, 행렬 $\bm E^1$과 $\bm C^1_k$ 또한 자동으로 순환 행렬이 되며, 특징 벡터 $\bm z \in \Re^d$에 대한 적용은 순환 합성곱 "$\circledast$"을 사용하여 구현될 수 있다.
구체적으로, 우리는 다음 명제를 갖는다.

\begin{proposition}[$\bm E^1$과 $\bm C^1_k$의 합성곱 구조]
행렬
\begin{equation}
    \bm E^1 = \alpha\big(\bm I + \alpha \circm(\bm Z^1) \circm(\bm Z^1)^\top \big)^{-1}
\end{equation}
은 순환 행렬이며 순환 합성곱을 나타낸다:
$$\bm E^1 \z = \bm e_1 \circledast \z,$$
여기서 $\bm e_1 \in \Re^d$은 $\bm E^1$의 첫 번째 열 벡터이고 "$\circledast$"는 다음과 같이 정의된 순환 합성곱이다.
\begin{equation*}
    (\bm e_1 \circledast \bm z)_{i} \doteq \sum_{j=0}^{d-1} e_1(j) x(i+ d-j \,\, \textsf{mod} \,\,d).
\end{equation*}
유사하게, $\bm Z^1$의 임의의 부분 집합과 관련된 행렬 $\bm C^1_k$ 또한 순환 합성곱이다.
\label{prop:circular-conv-1}
\end{proposition}

ReduNet의 첫 번째 계층 매개변수 $\bm E^1$과 $\bm{C}^1_k$가 순환 합성곱이 될 뿐만 아니라 다음 계층의 특징들도 순환 행렬로 남는다.
즉, \eqref{eqn:layer-approximate}의 증분 특징 변환이 $\z^1 \in \Re^d$의 모든 이동된 버전에 적용될 때, 다음과 같이 주어진다.
\begin{equation}
    \circm(\z^1) + \eta \cdot \bm E^1 \circm(\z^1) - \eta \cdot \bm \sigma \Big([\bm{C}_1^1 \circm(\z^1), \ldots, \bm{C}^1_K \circm(\z^1)] \Big),
\end{equation}
은 순환 행렬이다. 이는 첫 번째 계층에서 했던 것처럼 두 번째 계층 특징으로부터 순환 족을 구성할 필요가 없음을 의미한다.
다음과 같이 표기함으로써
\begin{equation}
\z^{2} \propto \z^{1} +\eta \cdot g(\z^{1}, \bm \theta^{1}) =  \z^1 + \eta \cdot \bm e_{1} \circledast \z^{1} -  \eta \cdot \bm \sigma\Big([\bm{c}_{1}^{1} \circledast \z^{1}, \dots, \bm{c}^{1}_{K} \circledast \z^{1}]\Big),
\label{eqn:approximate-convolution}
\end{equation}
다음 수준의 특징은 다음과 같이 쓸 수 있다.
$$\circm(\bm Z^2) = \big[ \circm( \bm z_{1}^1 + \eta g( \bm z_{1}^1, \bm \theta^1)), \dots, \circm( \bm z_{N}^1 + \eta g(\bm z_{N}^1, \bm \theta^1)) \big].$$
귀납적으로 계속하면, 이러한 $\circm(\bm Z^\ell)$을 기반으로 하는 모든 행렬 $\bm E^\ell$과 $\bm C^\ell_k$가 순환 행렬이며, 모든 특징도 마찬가지임을 알 수 있다.
데이터의 속성 덕분에, ReduNet은 {\em 이 구조를 명시적으로 선택할 필요 없이} 합성곱 네트워크의 형태를 띠게 되었다!

\paragraph{불변성과 희소성 사이의 근본적인 트레이드오프.} 
하지만 한 가지 문제가 있다: 일반적으로, 벡터 $\z$의 모든 원형 순열 집합은 전체 순위 행렬을 제공한다. 즉, 각 샘플(따라서 각 클래스)과 관련된 $d$개의 "증강된" 특징은 일반적으로 이미 전체 공간 $\Re^d$를 생성한다. 예를 들어, 델타 함수 $\delta(d)$의 모든 이동된 버전은 다른 모든 신호를 그들의 (밀집된) 가중 중첩으로 생성할 수 있다. MCR$^2$ 목표 \eqref{eqn:maximal-rate-reduction}는 클래스를 다른 부분 공간으로 구별할 수 없을 것이다.

한 가지 자연스러운 해결책은 원래 신호를 더 높은 차원의 공간으로 "리프팅"하여 데이터의 분리성을 향상시키는 것이다. 예를 들어, 여러 필터 $\bm k_1, \ldots, \bm k_C \in \Re^d$에 대한 응답을 취하는 것이다:
\begin{equation}
\bm z[c] = \bm k_c \circledast \bm x  =  \circm(\bm k_c) \bm x \in \Re^d, \quad c = 1, \ldots, C.
\label{eqn:lift-1d}
\end{equation} 
필터는 사전 설계된 불변성 촉진 필터일 수도 있고,\footnote{오디오와 같은 1D 신호의 경우, 기존의 단시간 푸리에 변환(STFT)을 고려할 수 있다; 2D 이미지의 경우, ScatteringNet \cite{scattering-net}에서와 같이 2D 웨이블릿을 고려할 수 있다.} 데이터로부터 적응적으로 학습될 수도 있으며,\footnote{학습된 필터의 경우, PCANet \cite{chan2015pcanet}에서와 같이 샘플의 주성분으로 필터를 학습하거나 합성곱 사전 학습 \cite{li2019multichannel,qu2019nonconvex}에서 학습할 수 있다.} 또는 우리 실험에서처럼 무작위로 선택될 수도 있다. 이 연산은 각 원래 신호 $\x \in \Re^d$를 $\bar{\z}  \doteq [\z[1], \ldots, \z[C]]^\top \in \Re^{C\times d}$로 표시되는 $C$-채널 특징으로 리프팅한다. 
그런 다음, $\bar{\z}$의 벡터 표현에 대해 ReduNet을 구성할 수 있다. 이는 다음과 같이 표시된다.
$\vec(\bar\z) \doteq [\z[1]^\top, \ldots, \z[C]^\top] \in \Re^{dC}$. 
관련된 순환 버전 $ \circm(\bar{\z})$과 모든 이동된 버전에 대한 데이터 공분산 행렬 $\bar{\bm \Sigma}(\bar\z)$는 다음과 같이 주어진다:
\begin{equation}
\begin{aligned}
 \circm(\bar{\z}) \doteq \left[\begin{matrix}
    \circm(\z[1])  \\ \vdots \\ \circm(\z[C]) \end{matrix} \right] \in \Re^{dC\times d}, 
    \quad  \bar{\bm \Sigma}(\bar\z) \doteq 
    \left[\begin{matrix}
    \circm(\z[1]) \\ \vdots \\ \circm(\z[C]) \end{matrix} \right]
    \left[\begin{matrix}\circm(\z[1])^\top,\ldots, \circm(\z[C])^\top\end{matrix} \right] \in \Re^{dC\times dC},
    \label{eqn:W-multichannel}
\end{aligned}
\end{equation}
여기서 $c \in [C]$인 $\circm(\z[c]) \in \Re^{d\times d}$는 특징 $\bar \z$의 $c$-번째 채널의 순환 버전이다. 그러면 $\circm(\bar\z)$의 열들은 $\Re^{dC}$에서 최대 $d$-차원의 적절한 부분 공간만을 생성할 것이다. 그러나 이 간단한 리프팅 연산(선형인 경우)은 아직 클래스를 분리 가능하게 만들기에 충분하지 않다 --- 다른 클래스와 관련된 특징들은 {\em 동일한} $d$-차원 부분 공간을 생성할 것이다. 이는 불변성과 선형 (부분 공간) 모델링 사이의 근본적인 충돌을 반영한다: {\em 임의로 이동되고 중첩된 신호가 동일한 클래스에 속하기를 바랄 수는 없다.}

\begin{figure}[t]
	\centerline{
\includegraphics[width=0.9\textwidth]{\toplevelprefix/chapters/chapter4/figs/sparse-representation.pdf}}
	\caption{각 입력 신호 $\bm x$(여기서는 이미지)는 사전 $\bm D$의 여러 커널 $\bm d_c$와의 희소 합성곱의 중첩으로 표현될 수 있다.}
	\label{fig:multi-channel-sparse-representation}
\end{figure}
이 충돌을 해결하는 한 가지 방법은 각 클래스 내의 추가적인 구조를 {\em 희소성}의 형태로 활용하는 것이다: 각 클래스 내의 신호는 일부 기본 원자(또는 모티프)의 임의의 선형 중첩으로 생성되는 것이 아니라, \Cref{fig:multi-channel-sparse-representation}에서 보여주듯이 그것들과 그것들의 이동된 버전의 {\em 희소한} 조합으로만 생성된다. 더 정확하게 말하면, 클래스 $k$와 관련된 원자 모음인 행렬 $\bm D_k = [\bm d_{k,1}, \ldots, \bm d_{k,c}]$를 사전이라고 하면, 이 클래스의 각 신호 $\x$는 다음과 같이 희소하게 생성된다:
\begin{equation}
    \x = \bm d_{k,1} \circledast z_1 + \ldots + \bm d_{k,c} \circledast z_c = \circm(\bm{D}_k)\z,
\end{equation}
어떤 희소 벡터 $\z$에 대해. 다른 클래스의 신호는 원자(또는 모티프)가 서로 비일관적인 다른 사전에 의해 생성된다. 비일관성 때문에, 한 클래스의 신호는 다른 클래스의 원자에 의해 희소하게 표현될 가능성이 낮다. 따라서 모든 신호는 다음과 같이 표현될 수 있다.
\begin{equation}
\x = \big[\circm(\bm{D}_1), \circm(\bm{D}_2), \ldots, \circm(\bm{D}_K)\big] \bar \z,
\end{equation}
여기서 $\bar \z$는 희소하다.\footnote{유사한 희소 표현 모델이 얼굴 인식과 같은 응용 분야에서 분류 목적으로 오랫동안 제안되고 사용되어 왔으며, 뛰어난 효과를 보여주었다 \cite{Wright:2009,wagner2012toward}. 최근에는 \cite{papyan2017convolutional}에 의해 심층 합성곱 네트워크의 구조를 해석하기 위한 프레임워크로 합성곱 희소 코딩 모델이 제안되었다.} 샘플 데이터로부터 가장 간결하고 최적의 희소화 사전을 학습하는 방법에 대한 방대한 문헌이 있으며, 예를 들어 \cite{li2019multichannel,qu2019nonconvex}가 있고, 이후 역 문제를 해결하고 관련된 희소 코드 $\z$ 또는 $\bar \z$를 계산한다. \cite{qu2020nonconvex,Qu2020Geometric}의 최근 연구는 넓은 조건 하에서 합성곱 사전 학습 문제가 효과적이고 효율적으로 해결될 수 있음을 보여주기까지 한다.

그럼에도 불구하고, 분류와 같은 작업의 경우, 우리는 반드시 각 개별 신호에 대한 정확한 최적 사전이나 정확한 희소 코드에 관심이 있는 것은 아니다. 우리는 주로 각 클래스에 대한 희소 코드 집합이 다른 클래스의 코드 집합과 적절하게 분리될 수 있는지에 관심이 있다. 희소 생성 모델의 가정 하에, 합성곱 커널 $\{\bm k_c\}_{c=1}^C$이 위 희소화 사전 $\bm D = [\bm D_1, \ldots, \bm D_K]$의 "전치" 또는 "역"과 잘 일치하면, 이를 {\em 분석 필터}라고도 하며 \cite{Cosparse-Nam,Analysis-Filter}, 한 클래스의 신호는 해당 필터의 작은 부분 집합에만 높은 응답을 보이고 다른 필터에는 낮은 응답을 보일 것이다(비일관성 가정 때문). 그럼에도 불구하고, 실제로는 종종 충분히 많은 수의, 예를 들어 $C$개의, 무작위 필터 $\{\bm k_c\}_{c=1}^C$가 추출된 $C$-채널 특징을 보장하기에 충분하다.
\begin{equation}
\big[\bm k_1 \circledast \x, \bm k_2 \circledast \x, \ldots, \bm k_C \circledast \x\big]^\top = \big[\circm(\bm k_1) \x, \ldots, \circm(\bm k_C) \x \big]^\top \in \Re^{C\times d}
\end{equation}
서로 다른 클래스에 대해 서로 다른 필터에 대한 응답 패턴이 다르므로 서로 다른 클래스를 분리할 수 있다 \cite{chan2015pcanet}.

따라서, 우리 프레임워크에서 채널 수(또는 네트워크의 너비)는 진정으로 {\em 통계적 자원}으로서의 역할을 하는 반면, 계층 수(네트워크의 깊이)는 {\em 계산적 자원}으로서의 역할을 한다. 압축 감지 이론은 데이터의 본질적인 저차원 구조(분리성 포함)를 보존하기 위해 얼마나 많은 측정이 필요한지를 정확하게 특성화한다 \cite{Wright-Ma-2021}. %

\begin{figure}[t]
	\centerline{
\includegraphics[width=0.95\textwidth]{\toplevelprefix/chapters/chapter4/figs/sparse-lifting.pdf}}
\caption{입력 신호 $\bm x$(여기서는 이미지)의 희소 코드 $\bar{\bm z}$를 여러 커널 $\bm k_c$와의 합성곱을 취한 후 희소화하여 추정한다.}
		\label{fig:multi-channel-sparse-lifting}
\end{figure}
다중 채널 응답 $\bar \z$는 희소해야 한다. 따라서 희소 코드 $\bar \z$를 근사하기 위해, 위 필터 출력에 대해 항목별 {\em 희소성 촉진 비선형 임계값}, 예를 들어 $\bm \tau(\cdot)$를 취하여 낮은(예: 절대값이 $\epsilon$ 미만) 또는 음수 응답을 0으로 설정할 수 있다:
\begin{equation}
\bar \z \doteq \bm \tau \left( \big[\circm(\bm k_1) \x, \ldots, \circm(\bm k_C) \x \big]^\top \right) \in \Re^{C \times d}.
\label{eqn:sparse-lifting}
\end{equation} 
\Cref{fig:multi-channel-sparse-lifting}은 기본 아이디어를 보여준다. 희소화 임계값 연산자의 설계에 대한 보다 체계적인 연구는 \cite{Analysis-Filter}를 참조할 수 있다. 그럼에도 불구하고, 여기서 우리는 코드가 충분히 분리될 수 있는 한 최상의 희소 코드를 얻는 데 그다지 관심이 없다. 따라서 비선형 연산자 $\bm \tau$는 소프트 임계값 또는 ReLU로 간단히 선택될 수 있다.
이러한 추정된 희소 특징 $\bar \z$는 $\mathbb{R}^{dC}$의 저차원 (비선형) 부분 다양체에 놓여 있다고 가정할 수 있으며, 이는 후속 ReduNet 계층에 의해 선형화되고 다른 클래스로부터 분리될 수 있다. 이는 나중에 \Cref{fig:learn-to-classify-diagram}에서 설명된다.

이러한 다중 채널 특징 $\bar\Z \doteq [\bar \z_1, \ldots, \bar \z_N] \in \Re^{C \times d \times N}$의 순환 버전, 즉 $\circm(\bar\Z) \doteq [ \circm(\bar\z_1), \dots, \circm(\bar\z_N)] \in \Re^{dC \times dN}$으로부터 구성된 ReduNet은 위에서 설명한 좋은 불변성 속성을 유지한다: 이제 $\bar{\bm E}$와 $\bar{\bm C}_k$로 표시되는 선형 연산자는 블록 순환 상태를 유지하며, {\em 다중 채널 1D 순환 합성곱}을 나타낸다.
구체적으로, 우리는 다음 결과를 얻는다.
\begin{proposition}[$\bar{\bm E}$와 $\bar{\bm C}_k$의 다중 채널 합성곱 구조]
행렬
\begin{equation}
\label{eq:def-E-bar}
\bar{\bm E} \doteq \alpha\left(\bm I + \alpha\, \circm(\bar\Z) \circm(\bar\Z)^\top \right) ^{-1}  
\end{equation}
은 블록 순환 행렬이다. 즉,
\begin{equation*}
    \bar{\bm E} = 
    \left[\begin{matrix}
        \bar{\bm E}_{1, 1} & \cdots & \bar{\bm E}_{1, C}\\
        \vdots & \ddots & \vdots \\
        \bar{\bm E}_{C, 1} & \cdots & \bar{\bm E}_{C, C}\\
    \end{matrix}\right] \in \Re^{dC \times dC},
\end{equation*}
여기서 각 $\bar{\bm E}_{c, c'}\in \Re^{d \times d}$는 순환 행렬이다. 더욱이, $\bar{\bm E}$는 다중 채널 순환 합성곱을 나타낸다. 즉, 임의의 다중 채널 신호 $\bar\z \in \Re^{C \times n}$에 대해
$$\bar{\bm E} \cdot \textsf{vec}(\bar\z) = \textsf{vec}( \bar{\bm e} \circledast \bar\z)$$ 
이 성립한다. 위에서 $\bar{\bm e} \in \Re^{C \times C \times d}$는 다중 채널 합성곱 커널이며, $\bar{\bm e}[c, c'] \in \Re^{d}$는 $\bar{\bm E}_{c, c'}$의 첫 번째 열 벡터이고, $\bar{\bm e} \circledast \bar\z \in \Re^{C \times d}$는 다음과 같이 정의된 다중 채널 순환 합성곱이다.
\begin{equation*}
    (\bar{\bm e} \circledast \bar\z)[c] \doteq \sum_{c'=1}^C \bar{\bm e}[c, c'] \circledast \bar{\z}[c'], \quad \forall c = 1, \ldots, C.
\end{equation*}
유사하게, $\bar{\bm Z}$의 임의의 부분 집합과 관련된 행렬 $\bar{\bm C}_k$ 또한 다중 채널 순환 합성곱이다.
\label{prop:multichannel-circular-conv-1}
\end{proposition}
\Cref{prop:multichannel-circular-conv-1}로부터, 이동 불변 ReduNet은 구성상 다중 채널 1D 신호를 위한 심층 합성곱 네트워크이다. 초기 리프팅 커널이 분리되더라도 \eqref{eqn:sparse-lifting}, $\bar{\bm E}$(그리고 유사하게 $\bar{\bm C_k}$)를 계산하기 위한 \eqref{eq:def-E-bar}의 행렬 역산은 모든 $C$ 채널 간에 "교차 토크"를 도입한다는 점에 유의하라. %
따라서, 이러한 다중 채널 합성곱은 일반적으로 한때 다중 채널 합성곱 신경망을 단순화하기 위해 제안되었던 Xception 망 \cite{Xception}과 달리 깊이별로 분리 가능하지 {\em 않다}.\footnote{데이터에 대한 어떤 추가 구조가 깊이별로 분리 가능한 합성곱으로 이어질지는 미해결 문제로 남아 있다.}

\begin{remark}[주파수 영역에서 계산 복잡도 줄이기]
\eqref{eq:def-E-bar}에서 $\bar{\bm E}$의 계산은 크기가 $dC \times dC$인 행렬을 역산해야 하며, 이는 일반적으로 복잡도가 $O(d^3C^3)$이다. 그럼에도 불구하고, 순환 행렬이 이산 푸리에 변환(DFT) 행렬에 의해 대각화될 수 있다는 사실을 이용하면 복잡도를 크게 줄일 수 있다. \cite{chan2021redunet}에서 보여주듯이, $\bar{\bm E}$와 $\bar{\bm C}_k \in \Re^{dC \times dC}$를 계산하기 위해, 우리는 주파수 영역에서 $C\times C$ 블록의 역행렬을 $d$번 계산하기만 하면 되므로 전체 복잡도는 $O(dC^3)$이 된다.
    
\end{remark}

\paragraph{전체 네트워크 아키텍처 및 비교.}
위의 유도 과정을 따라, 우리는 평행 이동에 불변하는 여러 클래스의 신호/이미지에 대해 선형 판별 표현(LDR)을 찾기 위해, 희소 코딩, 다중 채널 합성곱을 사용한 다중 계층 아키텍처, 다양한 비선형 활성화, 스펙트럼 컴퓨팅이 모두 목표를 효과적이고 효율적으로 달성하기 위한 {\em 필수적인} 구성 요소가 됨을 알 수 있다. \Cref{fig:learn-to-classify-diagram}은 입력 희소 코드에 대한 불변 레이트 감소를 통해 이러한 표현을 학습하는 전체 과정을 보여준다.

\begin{figure}[t]
    \centering
    \includegraphics[width=0.98\linewidth]{\toplevelprefix/chapters/chapter4/figs/learn_to_classify_diagram_updated.png}
    \caption{이동 불변성을 가진 다중 클래스 신호를 분류하기 위한 전체 과정: 다중 채널 리프팅, 희소 코딩, 그리고 불변 레이트 감소를 위한 다중 채널 합성곱 ReduNet. 이러한 구성 요소들은 이동 불변 다중 클래스 신호를 LDR로서 비일관성(선형) 부분 공간에 매핑하기 위해 {\em 필수적}이다. 대부분의 현대 심층 신경망의 아키텍처가 이 과정을 닮았다는 점에 유의하라. 이렇게 학습된 LDR은 분류와 같은 후속 작업을 용이하게 한다.}
    \label{fig:learn-to-classify-diagram}
\end{figure}

\begin{example}[숫자의 불변 분류]
다음으로 실제 10개 클래스의 MNIST 데이터셋에서 \textit{회전} 불변 특징을 학습하는 ReduNet의 경험적 성능을 제공한다.
이미지 $\bm{x}\in\mathbb{R}^{H\times W}$에 극좌표 그리드를 부과하며, 기하학적 중심이 2D 극좌표 그리드의 중심이 되도록 한다(\Cref{fig:samples-invariant-1d-mnist-diagram}에 도시됨).
각 반지름 $r_i$, $i \in [C]$에 대해, 각도 $\gamma_l =l\cdot({2\pi}/\Gamma)$, $l \in [\Gamma]$에 대해 $\Gamma$개의 픽셀을 샘플링할 수 있다.
그런 다음 데이터셋에서 샘플 이미지 $\bm{x}$가 주어지면, (샘플링된) 극좌표에서 이미지를 다중 채널 신호 $\bm{x}_p \in\R^{\Gamma\times C}$로 표현한다.
여기서 목표는 회전 불변 표현을 학습하는 것이다. 즉, $\{f(\bm{x}_p \circ \mathfrak{g}, \bm{\theta})\}_{\mathfrak{g} \in\mathbb{G}}$가 동일한 부분 공간에 있도록 $f(\cdot, \bm{\theta})$를 학습할 것으로 기대한다. 여기서 $\mathfrak{g}$는 극각에서의 순환 이동이다.
$N=100$개의 훈련 샘플(각 클래스에서 10개)을 사용하고, 극좌표 샘플링을 위해 $\Gamma=200, C=15$로 설정한다.
극좌표에서 위 샘플링을 수행함으로써, 데이터 행렬 $\bm{X}_p \in \mathbb{R}^{(\Gamma\cdot C) \times N}$를 얻을 수 있다.
ReduNet에 대해, 계층/반복 횟수 $L=40$, 정밀도 $\epsilon=0.1$, 스텝 크기 $\eta=0.5$로 설정한다. 첫 번째 계층 전에, 크기가 5인 20개의 무작위 가우시안 커널로 1D 순환 합성곱을 통해 입력 리프팅을 수행한다.

\begin{figure}[t]
    \begin{subfigure}[t]{0.3\textwidth}
        \centering
        \includegraphics[width=\textwidth]{\toplevelprefix/chapters/chapter4/figs/1d-rotation.pdf}
        \caption{$\bm{X}_{\text{rotation}}$}
    \end{subfigure}
    \hfill
    \begin{subfigure}[t]{0.3\textwidth}
        \centering
        \includegraphics[width=\textwidth]{\toplevelprefix/chapters/chapter4/figs/mnist1d_img0.pdf}
        \caption{$\bm{Z}_{\text{rotation}}$}
    \end{subfigure}
    \hfill
    \begin{subfigure}[t]{0.3\textwidth}
        \centering
        \includegraphics[width=\textwidth]{\toplevelprefix/chapters/chapter4/figs/mnist1d_img1.pdf}
        \caption{손실}
    \end{subfigure}
    \caption{\small MNIST 숫자들의 회전된 이미지 예시, 각각 18$^{\circ}$씩. (\textbf{왼쪽}) 극좌표 표현을 위한 다이어그램; (\textbf{오른쪽}) 숫자 `0'과 `1'의 회전된 이미지.}
    \label{fig:samples-invariant-1d-mnist-diagram}
\end{figure}


학습된 표현을 평가하기 위해, 각 훈련 샘플은 20개의 회전된 버전으로 증강되며, 각 버전은 보폭=10으로 이동된다. $m \times 20$개의 증강된 훈련 입력 $\bm{X}_{\text{rotation}}$ 사이의 코사인 유사도를 계산하고 그 결과는 \Cref{fig:redu-invariant-1d-mnist-diagram} (\textbf{a})에 표시된다.
모든 증강된 버전의 학습된 특징, 즉 $\bar{\bm{Z}}_{\text{rotation}}$ 사이의 코사인 유사도를 비교하고 그 결과를 \Cref{fig:redu-invariant-1d-mnist-diagram} (\textbf{b})에 요약한다.
보는 바와 같이, 이렇게 구성된 회전 불변 ReduNet은 10개의 다른 클래스로부터의 훈련 데이터(및 모든 회전된 버전)를 10개의 거의 직교하는 부분 공간으로 매핑할 수 있다. 즉, 학습된 부분 공간은 극각에서의 이동 변환에 진정으로 불변하다. 다음으로, 동일한 증강 절차를 따르는 또 다른 100개의 테스트 샘플을 무작위로 추출한다.
\Cref{fig:redu-invariant-1d-mnist-diagram} (\textbf{c})에서는 훈련 및 테스트 데이터셋에 대한 ReduNet의 $\ell$-번째 계층 표현에 대한 MCR$^{2}$ 손실을 시각화한다. 이러한 결과로부터, 구성된 ReduNet이 실제로 MCR$^{2}$ 손실을 최대화할 수 있을 뿐만 아니라 테스트 데이터에도 일반화됨을 알 수 있다.



\begin{figure}[t]
    \begin{subfigure}[t]{0.3\textwidth}
        \centering
        \includegraphics[width=\textwidth]{\toplevelprefix/chapters/chapter4/figs/mnist1d-heatmap-X_translate_train_all.pdf}
        \caption{$\bm{X}_{\text{rotation}}$}
    \end{subfigure}
    \hfill
    \begin{subfigure}[t]{0.3\textwidth}
        \centering
        \includegraphics[width=\textwidth]{\toplevelprefix/chapters/chapter4/figs/mnist1d-heatmap-Z_translate_train_all.pdf}
        \caption{$\bm{Z}_{\text{rotation}}$}
    \end{subfigure}
    \hfill
    \begin{subfigure}[t]{0.32\textwidth}
        \centering
        \includegraphics[width=\textwidth]{\toplevelprefix/chapters/chapter4/figs/mnist1d-loss-traintest.pdf}
        \caption{손실}
    \end{subfigure}
    \caption{\small (a)(b)는 회전 불변성에 대한 회전된 훈련 데이터 $\bm{X}_{\text{rotation}}$와 학습된 특징 $\bar{\bm{Z}}_{\text{rotation}}$ 사이의 코사인 유사도 히트맵이다. (d)는 계층에 걸친 훈련/검증 MCR$^2$ 손실을 시각화한다.}
    \label{fig:redu-invariant-1d-mnist-diagram}
\end{figure}

\end{example}

\section{펼쳐진 최적화를 통한 화이트박스 트랜스포머}\label{sec:chap4-white-box-transformer}
이전 섹션에서 보았듯이, 우리는 분류 문제를 사용하여 ResNet 및 CNN과 같은 인기 있는 심층 네트워크의 주요 아키텍처 특성에 대한 엄격한 해석을 제공했다. 즉, 이러한 네트워크의 각 계층은 레이트 감소(또는 정보 이득) 목표를 증가시키는 경사 단계를 모방하는 것으로 볼 수 있다. 이 관점은 또한 다소 놀라운 사실로 이어진다: ReduNet과 같은 심층 네트워크 계층의 매개변수와 연산자는 순전히 순방향 방식으로 계산될 수 있다.

ReduNet의 이론적, 개념적 중요성에도 불구하고, 몇 가지 요인이 매우 실용적이지 못하게 제한한다. 첫째, 위에서 논의했듯이, 순방향 방식으로 각 계층의 행렬 연산자를 계산하는 계산 비용이 매우 높을 수 있다. 둘째, 그렇게 계산된 연산자는 목표를 최적화하는 데 그다지 효과적이지 않을 수 있으며 수천 번의 반복(따라서 계층)이 필요할 수 있다. LISTA에 대한 \Cref{sec:LISTA}에서 보았듯이, 이러한 두 가지 문제는 해당 연산자를 최적화하고 역전파를 통해 학습 가능하게 함으로써 해결될 수 있다.\footnote{또는 순방향 및 역방향 최적화의 혼합을 통해 해결할 수도 있다.}

ReduNet이 유도된 지도 분류 설정도 다소 제한적이다. 실제 이미지에는 여러 객체가 포함될 수 있으므로 단일 클래스에 속하지 않을 수 있다. 따라서 이미지의 다른 영역이 다른 저차원 모델(가우시안 또는 부분 공간)에 속한다고 가정하는 것이 더 일반적일 것이다. 앞으로 보겠지만, 이러한 일반화는 이전 장에서 보았던 레이트 감소와 잡음 제거 연산을 통합하는 간단하고 일반적인 아키텍처로 이어진다. 더욱이, 그렇게 얻은 아키텍처는 인기 있는 트랜스포머 아키텍처와 유사하다.



\subsection{희소 레이트 감소를 위한 펼쳐진 최적화}



우리는 실제 신호와 관련된 일반적인 학습 설정을 고려한다. \(\X = \mat{\x_{1}, \dots, \x_{N}} \in \bR^{D \times N}\)을 데이터 소스를 나타내는 확률 변수라고 하자. 비전 작업에서 각 \(\x_{i} \in \bR^{D}\)는 일반적으로 이미지 패치에 해당하는 \textit{토큰}으로 해석된다. 언어 작업에서 각 \(\x_{i} \in \bR^{D}\)는 단어나 하위 단어와 같은 이산 토큰의 연속 벡터 표현인 \textit{토큰 임베딩}으로 해석된다.\footnote{용어의 약간의 남용과 함께, 이 장에서는 편의상 이산 토큰과 관련 임베딩을 모두 단순히 토큰이라고 부른다.}  %
\(\x_{i}\)는 임의의 상관 구조를 가질 수 있다. 우리는 \(\Z = \mat{\z_{1}, \dots, \z_{N}} \in \bR^{d \times N}\)을 우리의 표현을 정의하는 확률 변수를 나타내는 데 사용하며, 여기서 \(\z_{i} \in \bR^{d}\)는 해당 토큰 \(\x_i \in \bR^{D}\)의 표현이다.%

\begin{remark}
    트랜스포머에서 각 입력 샘플은 일반적으로 {\em 토큰}의 시퀀스로 변환된다. 토큰은 원시 입력에서 파생된 정보의 기본 단위이다: 자연어 처리에서 토큰은 일반적으로 단어 또는 하위 단어이다; 컴퓨터 비전에서는 이미지 패치에 해당한다; 다른 양식에서는 시간 단계, 공간 위치 또는 기타 도메인별 단위를 나타낼 수 있다. {\em 토큰 임베딩}은 트랜스포머에 대한 입력으로 사용되는 토큰의 연속 벡터 표현이다. 각 토큰을 고차원 공간의 한 점에 매핑하여 모델이 수치 계산을 사용하여 기호 입력을 처리할 수 있도록 한다.
    {\em 토큰 표현}은 일반적으로 트랜스포머의 중간 또는 최종 계층에 의해 생성되는 토큰의 의미론적 또는 구조적 정보를 인코딩하는 벡터이다. 이러한 표현은 분류, 생성 또는 회귀와 같은 다운스트림 작업에 유용한 입력의 의미 있는 특징을 포착하도록 설계되었다. 구현에서 이러한 개념에 대한 자세한 내용은 \Cref{sec:contrastive_learning}를 참조하라.
\end{remark}


 
\paragraph{구조화되고 간결한 표현 학습을 위한 목표.}
레이트 감소 프레임워크 \Cref{sec:chap4-white-box-model-via-unrolling}에 따라, 우리는 표현 학습의 목표가 잠재적으로 비선형적이고 다중 모드 분포를 갖는 입력 토큰 \(\{\bm x_i\}_{i=1}^N \subset \R^D\)을 (조각별) \textit{선형화되고 간결한} 토큰 표현 \(\{\bm z_i\}_{i=1}^N \subset \bR^{d}\)으로 변환하는 특징 매핑 \(f \colon \X \in \bR^{D \times N} \to \Z\in \bR^{d \times N}\)를 찾는 것이라고 주장한다. 토큰 표현 \(\{\z_{i}\}_{i = 1}^{N}\)의 결합 분포는 정교할 수 있지만(그리고 작업에 따라 다를 수 있지만), 우리는 개별 토큰 표현의 목표 주변 분포가 간결한 코딩에 적합하도록 고도로 압축되고 구조화되어야 한다고 요구하는 것이 합리적이고 실용적이라고 주장한다. 특히, 우리는 분포가 \textit{저차원(예: \(K\)개) 가우시안 분포의 혼합}이 되도록 요구하며, 여기서 \(k\)-번째 가우시안은 평균 \(\Zero \in \bR^{d}\), 공분산 \(\vSigma_{k} \succeq \Zero \in \bR^{d \times d}\)를 가지며, 직교 정규 기저 \(\vU_{k} \in \bR^{d \times p}\)에 의해 생성되는 지지(support)를 갖는다.
우리는 모든 가우시안의 기저 집합을 $\vU_{[K]} = \{\vU_k\}_{k=1}^K$로 나타낸다. 따라서 최종 토큰 표현에 대한 \textit{정보 이득} \cite{ma2022principles}을 최대화하기 위해, 우리는 그들의 레이트 감소를 최대화하기를 원한다(\Cref{subsec:MCR2} 참조). 즉,
\begin{align}\label{eq:rate reduction}
    \mathrm{max}_{\bm Z \in \R^{d\times N}}\ \Delta R_{\epsilon}(\Z \mid \vU_{[K]}) \doteq R_{\epsilon}(\Z) - R^c_{\epsilon}(\Z \mid \vU_{[K]}).
\end{align}
여기서 첫 번째 항 $R_{\epsilon}$은 전체 토큰 표현 집합에 대한 손실 압축 코딩 레이트의 추정치이다. 더 구체적으로, 토큰 표현 $\{\bm z_i\}_{i=1}^N$을 평균이 0인 단일 가우시안 분포에서 추출된 i.i.d. 샘플로 본다면, 양자화 정밀도 $\epsilon > 0$에 대한 손실 압축 코딩 레이트는 다음과 같이 주어진다.
\begin{equation}\label{eq:coding_rate}
    R_{\epsilon}(\Z) \doteq \frac{1}{2}\textrm{logdet}\left(\I + \frac{d}{N\epsilon^{2}}\Z^{\top}\Z\right) = \frac{1}{2}\textrm{logdet}\left(\I + \frac{d}{N\epsilon^{2}}\Z\Z^{\top}\right).
\end{equation}
두 번째 항 $R_{\epsilon}^c$는 코드북 $\bm U_{[K]}$ 하에서의 손실 압축 코딩 레이트의 추정치이며, 다음과 같이 주어진다.
\begin{equation}
\begin{aligned}
    R_{\epsilon}^c(\Z \mid \vU_{[K]}) &\doteq \sum_{k=1}^{K}R_{\epsilon}(\vU_k^{\top} \Z) = \frac{1}{2}\sum_{k=1}^{K}\log\det\left(\I +
    \frac{p}{N\epsilon^{2}}(\vU_k^{\top}\Z)^{\top}(\vU_k^{\top}\Z)\right).
\end{aligned}
\label{eq:def-mcr-Rc}
\end{equation}

\begin{remark}
    코딩 레이트에 대한 표현식 \eqref{eq:def-mcr-Rc}는 원래 레이트 감소 목표 \eqref{eq:MCRc}에서 사용된 코딩 레이트 \(R_\epsilon^{c}\)의 일반화로 볼 수 있다. 특히, 원래 목표는 특정 데이터 실현 \(\vX\)에 특정한 알려진 소속 레이블 집합 \(\{\vPi_{k}\}\)에 대해 정의된다. 반면, 현재 목표는 특정 실현과 무관하지만 토큰 표현의 분포를 지원하는 것으로 가정되는 부분 공간 \(\vU_{[K]}\)에 대해 정의된다. 토큰 표현 $\bm z_i$이 부분 공간 $\bm U_k$에 속하고 이러한 부분 공간들이 서로 거의 직교한다고 가정하자. 즉, 모든 $k \neq l$에 대해 $\bm U_k^\top \bm U_l \approx \bm 0$이다. 그러면 투영 $\bm U_k\bm U_k^\top \bm z_i  = \bm z_i$와 모든 $l \neq k$에 대해 $\bm U_l\bm U_l^\top \bm z_i \approx \bm 0$임을 확인할 수 있다. 이러한 직교 투영은 각 토큰 표현이 속한 부분 공간을 식별하는 암묵적인 소속 레이블 역할을 효과적으로 수행한다.
\end{remark}



\begin{figure}[t!]
     \centering
         \includegraphics[width=0.95\textwidth]{\toplevelprefix/chapters/chapter4/figs/coding-transform.png}
     \caption{ \small\textbf{레이트 감소와 희소성을 통한 세 가지 표현 집합의 비교.} 각 $S_i$는 하나의 선형 부분 공간을 나타내며, 파란색 공의 수는 코딩 레이트의 차이 $\Delta R_{\epsilon}(\vZ \mid \vU_{[K]}) = R_\epsilon(\vZ) - R^c_\epsilon(\vZ \mid \vU_{[K]})$를 나타낸다.
     }
        \label{fig:sparse-rate-reduction-diagram}
\end{figure}


\paragraph{희소 레이트 감소.} 레이트 감소 목표 \eqref{eq:rate reduction}는 표현과 부분 공간의 임의의 결합 회전에 불변하다는 점에 유의하라. 특히, 레이트 감소 목표를 최적화하는 것이 자연스럽게 축에 정렬된(즉, \textit{희소한}) 표현으로 이어지지 않을 수 있다. {예를 들어, \Cref{fig:sparse-rate-reduction-diagram}의 세 가지 학습된 표현 집합을 고려해 보자. 코딩 레이트 감소는 (a)에서 (b)로 증가하지만, 회전에 불변하기 때문에 (b)에서 (c)로 동일하게 유지된다.} 따라서, 우리는 표현(과 그들의 지지 부분 공간)을 변환하여 표현 $\vZ$이 결과적인 표현 공간의 표준 좌표에 대해 결국 희소해지도록 하고 싶다\footnote{구체적으로, 0이 아닌 항목이 거의 없음.}. {이는 \Cref{fig:sparse-rate-reduction-diagram}(c)에서와 같다.} 따라서, 최종 표현이 더 간결한 코딩에 적합하도록 보장하기 위해, 우리는 표현(과 그들의 지지 부분 공간)을 변환하여 결과적인 표현 공간의 표준 좌표에 대해 \textit{희소}해지도록 하고 싶다.\footnote{즉, 0이 아닌 항목이 가장 적음.} 계산적으로, 우리는 위 두 목표를 최적화를 위한 통일된 목표로 결합할 수 있다:
\begin{equation}
   \max_{f \in \mathcal{F}}\ [\Delta R_{\epsilon}(\Z \mid \vU_{[K]}) - \lambda \|\bm{Z}\|_0] \qquad \text{s.t.}\ \Z = f(\X),
   \label{eq:sparse-rr}
\end{equation}
여기서 $\mathcal{F}$는 일반적인 함수 클래스를 나타내고 $\ell_0$ 노름 $\|\Z\|_0$는 최종 토큰 표현 \(\Z = f(\X)\)의 희소성을 촉진한다.%


실제로, $\ell_0$ 노름은 계산 추적성을 향상시키고 볼록 최적화 기법을 가능하게 하기 위해 종종 $\ell_1$ 노름으로 완화된다 \cite{Wright-Ma-2022}. 이에 동기를 부여받아, 우리는 문제 \eqref{eq:sparse-rr}을 그에 따라 완화하여, 원래의 희소성 목표에 충실하면서도 다음과 같이 효율적인 알고리즘에 더 적합한 공식을 도출한다:
\begin{equation}
\begin{aligned}
   \max_{f \in \mathcal{F}}\ [\Delta R_{\epsilon}(\Z \mid \vU_{[K]}) - \lambda \|\bm{Z}\|_1]  \qquad \text{s.t.}\ \Z = f(\X),
   \label{eq:sparse-rr-1}
\end{aligned}
\end{equation}
용어의 약간의 남용과 함께, 우리는 종종 이 목적 함수를 \textit{희소 레이트 감소}라고도 부른다.

\paragraph{펼쳐진 최적화를 통한 화이트박스 네트워크 아키텍처.}


언급하기는 쉽지만, 위 목표의 각 항은 계산적으로 최적화하기 어렵다 \cite{Wright-Ma-2022}. 따라서 전역 변환 $f$를 실현하여 
\eqref{eq:sparse-rr}을 최적화하기 위해, 표현 분포를 원하는 간결한 모델 분포로 밀어 넣는 여러 개의(예: $L$개) 간단한 \textit{점진적이고 국소적인} 연산 $f^\ell$의 연결을 통해 근사 접근법을 채택하는 것이 자연스럽다:
\begin{equation}
f\colon \X = \bm Z^0 \xrightarrow{\hspace{1mm} f^0 \hspace{1mm}} \Z^1 \rightarrow \cdots \rightarrow \Z^\ell \xrightarrow{\hspace{1mm} f^{\ell} \hspace{1mm}} \Z^{\ell+1} \rightarrow  \cdots \xrightarrow{\hspace{1mm} f^{L-1}} \Z^L = \Z,
\label{eq:incremental}
\end{equation}
여기서 $f^0: \bR^{D} \rightarrow \bR^{d}$는 각 입력 토큰 $\x_{i} \in \bR^{D}$을 초기 토큰 표현 $\z_{i}^{1} \in \bR^{d}$로 변환하는 전처리 매핑이다.
각 점진적 \textit{순방향 매핑} $\Z^{\ell + 1} = f^\ell(\Z^\ell)$ 또는 "계층"은 입력 $\Z^\ell$의 분포에 따라 위의 희소 레이트 감소 목표 \eqref{eq:sparse-rr}를 \textit{최적화}하기 위해 토큰 분포를 변환한다.

\begin{remark}
    ReduNet(\Cref{sec:chap4-white-box-model-via-unrolling} 참조)과 같은 다른 펼쳐진 최적화 접근 방식과 달리, 우리는 각 계층에서 $\Z^\ell$의 분포를 명시적으로 모델링한다. 예를 들어, 선형 부분 공간의 혼합 또는 사전에서 희소하게 생성된 것으로 모델링한다. 모델 매개변수는 데이터로부터 학습된다(예: 종단 간 훈련을 통한 \textit{역전파}). 순방향 "최적화"와 역방향
"학습" 사이의 이러한 분리는 각 계층의 수학적 역할을 명확히 한다. 즉, 각 계층은 입력 분포를 변환하는 연산자이며, 입력 분포는 차례로 계층의 매개변수에 의해 모델링되고 (이후 학습됨)된다.
\end{remark}

이제, 우리는 문제 \eqref{eq:sparse-rr-1}을 해결하기 위해 펼쳐진 최적화 관점을 통해 이러한 점진적이고 국소적인 연산을 어떻게 유도하는지 보여준다. 문제
\eqref{eq:sparse-rr-1}을 최적화하기 위해 점진적인 접근 방식을 사용하기로 결정하면, 최적화를 달성하기 위한 다양한 선택이 가능하다. $\Z^\ell$에 대한 모델(예: 부분 공간의 혼합 $\vU_{[K]}$)이 주어지면, 우리는 강력한 개념적 기반을 가진 2단계 \textit{교대 최소화} 방법을 선택한다. 첫째, 코딩 레이트 항 $R^c_\epsilon (\vZ \mid \vU_{[K]}^\ell)$을 최소화하기 위해 경사 하강법을 통해 토큰 $\vZ^{\ell}$을 \textit{압축}한다. 구체적으로, 다음과 같이 학습률 $\kappa$로 $R^c_\epsilon$에 대한 경사 단계를 취한다:
\begin{align}\label{eq:Z l+1/2}
    \bm Z^{\ell+1/2} = \bm Z^\ell - \kappa \nabla_{\bm Z} R^c_\epsilon (\vZ \mid \vU_{[K]}^\ell). 
\end{align}
다음으로, 압축된 토큰을 \textit{희소화}하여, 남은 항 $\lambda \norm{\vZ}_{1} - R_{\epsilon}(\vZ)$을 최소화하기 위해 적절히 완화된 근접 경사 단계를 통해 \(\vZ^{\ell + 1}\)을 생성한다. 나중에 자세히 논의하겠지만, 우리는 사전 $\bm D^\ell$에 대한 희소 표현 문제를 해결하여 이러한 $\bm Z^{\ell+1}$을 찾을 수 있다:
\begin{equation}\label{eq:Z l+1}
  \vZ^{\ell+1} = \argmin_{{\vZ}}  \bigg\{\lambda \norm{\vZ}_1 + \frac{1}{2}\norm{\vZ^{\ell + 1/2} - \vD^{\ell} {\vZ}}_F^2\bigg\}.
\end{equation}
다음에서, 우리는 위의 두 단계 각각에 대한 기술적 세부 사항을 제공하고 구현을 위한 효율적인 업데이트를 유도한다.




\paragraph{토큰 표현의 코딩 레이트에 대한 경사 하강법으로서의 자기-어텐션.} 첫 번째 단계 \eqref{eq:Z l+1/2}의 경우, 코딩 레이트의 그래디언트 \(\nabla_{\bm Z} R^c_\epsilon\)는 계산 비용이 많이 든다. 이는 \(K\)개의 기저 \(\vU_{k}^{\ell}\)를 갖는 \(K\)개의 부분 공간 각각에 대해 하나씩, 총 \(K\)개의 개별 행렬 역행렬을 포함하기 때문이다:
\begin{equation}
    \nabla_{\bm Z} R_{\epsilon}^c(\vZ \mid \vU_{[K]})
    = \frac{p}{N\epsilon^2}\sum_{k=1}^K \vU_k\vU_k^\top\vZ\Big(\I +
    \frac{p}{N\epsilon^2}(\vU_k^\top\vZ)^\top(\vU_k^\top\vZ)\Big)^{-1}.
    \label{eq:rate-gradient}
\end{equation}
이제, 우리는 이 그래디언트가 소위 다중 헤드 부분 공간 자기-어텐션(MSSA) 연산자를 사용하여 자연스럽게 근사될 수 있음을 보여준다. 이 연산자는 \(K\)개의 헤드를 가진 다중 헤드 자기-어텐션 연산자\citep{vaswani2017attention}와 유사한 함수 형태를 가진다(즉, 각 부분 공간에 대해 하나씩, 각 행렬 역행렬에서 비롯됨). 여기서, 우리는 1차 노이만 급수(\Cref{ex:neumannn} 참조)를 사용하여 그래디언트 \eqref{eq:rate-gradient}를 근사한다:
\begin{align}\label{eq:gd_rc_neumann}
    \nabla_{\bm Z} R_{\epsilon}^{c}(\vZ \mid \vU_{[K]}) 
    &\approx \frac{p}{N\epsilon^2} \sum_{k = 1}^{K}\vU_{k}\vU_{k}^\top\vZ\left(\vI - \frac{p}{N\epsilon^2} (\vU_{k}^\top\vZ)^\top(\vU_{k}^\top\vZ)\right) \notag \\
    &{= \frac{p}{N\epsilon^2} \left(\sum_{k = 1}^{K} \vU_{k}\vU_{k}^\top\right)\vZ -  \left( \frac{p}{N\epsilon^2}\right)^2\sum_{k = 1}^{K} \vU_{k}(\vU_{k}^\top\vZ)(\vU_{k}^\top\vZ)^\top(\vU_{k}^\top\vZ)}.
\end{align}
이 근사에서, 우리는 투영된 특징 간의 자기 상관관계를 통해 투영된 토큰 표현 $\{\bm U_k^\top\bm z_i\}$ 간의 유사성을 계산하고, 이를 소프트맥스, 즉 $\softmax{(\vU_{k}^\top\vZ)^\top(\vU_{k}^\top\vZ)}$를 사용하여 소속 분포로 변환한다.
부분 공간들의 합집합 $\bm U_{[K]}$이 전체 공간을 생성한다고 가정하자. 그러면 $\sum_{k = 1}^{K} \vU_{k}\vU_{k}^\top = \bm I$이다. 따라서, \eqref{eq:gd_rc_neumann}는 다음과 같이 된다.
\begin{align}\label{eq:grad Rc}
    \nabla_{\bm Z} R_{\epsilon}^{c}(\vZ \mid \vU_{[K]}) 
     \approx  \frac{p}{N\epsilon^2} \vZ -  \left( \frac{p}{N\epsilon^2}\right)^2 \MSSA\left(\vZ^{\ell} \mid \vU_{[K]}^{\ell}\right),
\end{align}
여기서 MSSA는 SSA 연산자를 통해 다음과 같이 정의된다:
\begin{align}
    & \mathrm{SSA}\left(\vZ \mid \vU_{k}\right) 
    \doteq (\vU_{k}^\top \vZ)\mathrm{softmax}\left((\vU_{k}^\top\vZ)^\top(\vU_{k}^\top\vZ)\right), \ \forall k \in [K], \label{eq:SSA} \\
    & \mathrm{MSSA}\left(\vZ \mid \vU_{[K]}\right) 
    \doteq \frac{p}{N\epsilon^2} \mat{\vU_{1}, \dots, \vU_{K}}\mat{\mathrm{SSA}({\vZ \mid \vU_{1}}) \\ \vdots \\ \mathrm{SSA}({\vZ \mid \vU_{K}})}.\label{eq:Multi-Head-SSA}
\end{align}  
\eqref{eq:grad Rc}를 \eqref{eq:Z l+1/2}에 대입하면 다음과 같이 자연스럽게 근사될 수 있다.
\begin{equation}
    \vZ^{\ell + 1/2} = \left(1 - \frac{\kappa p}{N\epsilon^2}\right) \vZ^{\ell} + \frac{\kappa p}{N\epsilon^2} \mathrm{MSSA}\left(\vZ^{\ell}\ \middle|\ \vU_{[K]}^{\ell}\right).  \label{eq:gd-mcr-parts} 
\end{equation}



\begin{remark}
    \eqref{eq:SSA}의 SSA 연산자는 일반적인 트랜스포머의 \textit{어텐션 연산자}와 유사하지만\citep{vaswani2017attention}, 여기서는 값, 키, 쿼리의 선형 연산자가 모두 부분 공간 기저와 \textit{동일}하게 설정된다는 점이 다르다. 즉, $\vV_{k} = \vK_{k} = \bm{Q}_{k} = \vU_k^*$이다. 따라서, 우리는 $\mathrm{SSA}({\spcdot\mid\vU_k}): \bR^{d\times n} \rightarrow \bR^{p\times n}$를 \textbf{S}ubspace \textbf{S}elf-\textbf{A}ttention (SSA) 연산자라고 명명한다. 그러면, \eqref{eq:Multi-Head-SSA}의 전체 MSSA 연산자는 공식적으로 \(\mathrm{MSSA}({\spcdot \mid \vU_{[K]}}) \colon \bR^{d \times n} \to \bR^{d \times n}\)로 정의되며, \textbf{M}ulti-Head \textbf{S}ubspace \textbf{S}elf-\textbf{A}ttention (MSSA) 연산자라고 불린다. 이는 기존 트랜스포머 네트워크의 인기 있는 다중 헤드 자기-어텐션 연산자와 개념적으로 유사하게 모델 종속 가중치를 사용하여 어텐션 헤드 출력을 평균내어 집계한다. 전체 경사 단계 \eqref{eq:gd-mcr-parts}는 트랜스포머에서 스킵 연결로 구현된 다중 헤드 자기-어텐션과 유사하다.
\end{remark}




\paragraph{토큰 표현의 희소 코딩을 위한 근접 경사 하강법으로서의 MLP.} 교대 최소화의 두 번째 단계에서는 $\lambda \norm{\vZ}_{1} - R_{\epsilon}(\vZ)$를 최소화해야 한다. 그래디언트 \(\nabla R_{\epsilon}(\vZ)\)는 행렬 역행렬을 포함하므로, 이 문제를 최적화하기 위한 순진한 근접 경사법(\Cref{subsec:pgd} 참조)은 대규모 문제에서 다루기 어려워진다. %
따라서 우리는 표현의 다양성과 희소화 사이의 절충을 위해 다른 단순화 접근법을 취한다: (완전한) 비일관적이거나 직교하는 사전 $\vD^{\ell} \in \bR^{d \times d}$을 가정하고, 중간 반복 \(\vZ^{\ell + 1/2}\)을 \(\vD^{\ell}\)에 대해 희소화하도록 요청한다. 즉, $\vZ^{\ell + 1/2} \approx \vD^{\ell} \vZ^{\ell + 1}$이며, 여기서 $\vZ^{\ell + 1}$은 더 희소하다; 즉, \(\vZ^{\ell + 1/2}\)의 \textit{희소 인코딩}이다. 사전 \(\vD^{\ell}\)은 모든 토큰을 동시에 희소화하는 데 사용된다.
비일관성 가정에 의해, $(\vD^{\ell})^\top(\vD^{\ell}) \approx \vI$이다. 따라서 \eqref{eq:coding_rate}로부터 다음을 얻는다.
\begin{equation}
    R_{\epsilon}(\vZ^{\ell + 1/2}) \approx R_{\epsilon}(\vD^{\ell}\vZ^{\ell + 1}) \approx R_{\epsilon}(\vZ^{\ell + 1}).
\end{equation}  
$\lambda \norm{\vZ}_{1} - R_{\epsilon}(\vZ)$를 풀기 위해, 우리는 다음 문제를 최적화한다.
\begin{align*}
       \vZ^{\ell + 1} \approx \argmin_{\vZ}  \norm{\vZ}_1 \quad \mbox{subject to} \quad \vZ^{\ell + 1/2} = \vD^{\ell}\vZ.
\end{align*}
{위의 희소 표현 프로그램은 일반적으로 LASSO \citep{Wright-Ma-2022}로 알려진 제약 없는 볼록 프로그램으로 완화하여 해결된다:}
\begin{equation}
    \vZ^{\ell + 1} \approx \argmin_{\vZ} \left[\lambda \norm{\vZ}_1 + \frac{1}{2}\norm{\vZ^{\ell + 1/2} - \vD^{\ell} \vZ}_F^2 \right].
\end{equation}
우리 구현에서는 $\vZ^{\ell + 1}$에 음수가 아닌 제약 조건을 추가하고, 해당 음수 미포함 LASSO를 푼다:
\begin{equation}
    \vZ^{\ell + 1} \approx \argmin_{\vZ \geq \bm 0} \left[\lambda\norm{\vZ}_1 + \frac{1}{2}\norm{\vZ^{\ell + 1/2} - \vD^{\ell} \vZ}_{F}^{2}\right].
    \label{eq:sparse-nonnegative}
\end{equation}
그런 다음, ISTA 단계 \citep{beck2009fast}로 알려진 펼쳐진 {\em 근접 경사 하강법} 단계를 수행하여 \Cref{eq:sparse-nonnegative}를 점진적으로 최적화하여 다음과 같은 업데이트를 제공한다:
\begin{align}\label{eq:ista-block}
    \vZ^{\ell + 1} 
    &= \mathrm{ISTA}({\vZ^{\ell + 1/2} \mid \vD^{\ell}}), \\
    \text{where} \quad \mathrm{ISTA}({\vZ \mid \vD}) 
    &\doteq \operatorname{ReLU}(\vZ - \eta \vD^\top(\vD\vZ - \vZ) - \eta \lambda \bm{1}).
\end{align}






\begin{figure}
     \centering
     \includegraphics[width=0.75\textwidth]{\toplevelprefix/chapters/chapter4/figs/crate_encoder_architecture.pdf}
    \caption{\small \textbf{CRATE 인코더 아키텍처의 한 계층.} 전체 아키텍처는 단순히 이러한 계층들을 연결한 것이며, 일부 초기 토크나이저, 전처리 헤드, 그리고 최종 작업별 헤드(즉, 분류 헤드)가 있다.}
    \label{fig:crate_backbone}
\end{figure}



\subsection{전체 화이트박스 트랜스포머 아키텍처: CRATE}

이제 위의 업데이트를 펼쳐서 코딩 RATE 트랜스포머(\textsc{crate})라는 화이트박스 트랜스포머 아키텍처를 설계한다. 위의 두 단계 \eqref{eq:gd-mcr-parts}와 \eqref{eq:ista-block}을 결합하여:
\begin{enumerate}[leftmargin=0.7cm]
    \item 샘플 내 토큰을 부분 공간 혼합 구조로 국소적으로 압축하여 다중 헤드 부분 공간 자기-어텐션 블록(\texttt{MSSA})을 도출한다.
    \item 희소 코딩을 통해 모든 샘플에 걸쳐 토큰 집합을 전역적으로 희소화하여 희소화 블록(\texttt{ISTA})을 도출한다.
\end{enumerate}
우리는 \Cref{fig:crate_backbone}에 설명된 다음과 같은 레이트 감소 기반 트랜스포머 계층을 얻을 수 있다.
\begin{equation}
    \vZ^{\ell+1/2} \doteq \vZ^{\ell} + \texttt{MSSA}(\vZ^{\ell} \mid \vU_{[K]}^{\ell}), 
    \qquad 
    \vZ^{\ell+1}\doteq \texttt{ISTA}(\vZ^{\ell+1/2} \mid \bm D^\ell).
\end{equation}
우리의 표현을 점진적으로 구성하는 \eqref{eq:incremental}에 따라 여러 개의 이러한 계층을 구성함으로써, 우리는 데이터 토큰을 간결하고 희소한 비일관성 부분 공간의 합집합으로 변환하는 화이트박스 트랜스포머 아키텍처를 얻는다. 여기서 $f^{\pre}: \bR^{D \times N} \rightarrow \bR^{d \times N}$는 입력 토큰 $\vX \in \bR^{D \times N}$을 첫 번째 계층 표현 $\vZ^{1} \in \bR^{d \times N}$으로 변환하는 전처리 매핑이다. 이 아키텍처의 전체 흐름은 \Cref{fig:crate-diagram}에 나와 있다.

\begin{figure}[t!]
     \centering
         \includegraphics[width=\textwidth]{\toplevelprefix/chapters/chapter4/figs/CRATE_fig1_patches.pdf}
     \vspace{-0.1in}
     \caption{
     \textbf{\textsc{crate} 화이트박스 심층 네트워크 설계의 '주 루프'.}
     입력 데이터를 토큰 시퀀스 $\vZ^0$으로 인코딩한 후, \textsc{crate}는 분포에 대한 국소 모델에 대해 연속적으로 {\textit{\textbf{압축}}}하여 $\vZ^{\ell+1/2}$을 생성하고, 전역 사전에 대해 {\textit{\textbf{희소화}}}하여 $\vZ^{\ell+1}$을 생성함으로써 데이터를 저차원 부분 공간의 표준 구성으로 변환하는 심층 네트워크를 구성한다.
     이러한 블록을 반복적으로 쌓고 역전파를 통해 모델 매개변수를 훈련하면 데이터의 강력하고 해석 가능한 표현을 얻을 수 있다.
     }
        \label{fig:crate-diagram}
\end{figure}


\begin{remark}[\textbf{순방향 전파와 역전파의 역할}]\label{sub:forward_backward}
    ReduNet \cite{chan2021redunet}과 같은 다른 펼쳐진 최적화 접근 방식과 달리, 우리는 각 계층에서 각 $\vZ^\ell$과 $\vZ^{\ell + 1/2}$의 분포를 명시적으로 모델링한다. 이는 선형 부분 공간의 혼합 또는 사전에서 희소하게 생성된 것으로 이루어진다. 우리는 각 계층 \(\ell\)에서 학습된 부분 공간 기저 \(\vU_{[K]}^{\ell}\)와 학습된 사전 \(\vD^{\ell}\)이 함께 각 계층 \(\ell\)의 중간 표현을 인코딩하고 변환하는 \textit{코드북} 또는 \textit{분석 필터} 역할을 한다는 해석을 도입했다. 계층 \(\ell\)에 대한 입력 분포는 먼저 \(\vU_{[K]}^{\ell}\)에 의해 모델링된 다음 \(\vD^{\ell}\)에 의해 변환되므로, 각 계층에 대한 입력 분포는 다르다. 따라서 가장 간결한 인코딩을 얻으려면 각 계층에 별도의 코드북이 필요하다. 지금까지 완벽하게 알려져 있다고 가정했던 이러한 코드북의 매개변수(즉, 부분 공간 기저와 사전)는 실제로 데이터로부터 학습된다(예: 종단 간 훈련 내에서 \textit{역전파}를 통해).
    
    따라서, 우리의 방법론은 이렇게 파생된 화이트박스 심층 신경망에 대해 순방향 "최적화"와 역방향 "학습" 사이의 명확한 개념적 분리를 특징으로 한다. 즉, 순방향 전파에서는 각 계층을 입력 분포에 대한 학습된 모델(즉, 코드북)을 조건으로 하여 이 분포를 더 간결한 표현으로 변환하는 연산자로 해석한다. 역전파에서는 각 계층의 입력 분포에 대한 이 모델의 코드북이 특정 (지도) 입출력 관계에 더 잘 맞도록 업데이트된다. 이는 \Cref{fig:forward-backward}에 설명되어 있다. 이 개념적 해석은 특정 작업 및 손실에 대한 모델 표현의 특정 불가지론을 의미한다. 특히, 많은 유형의 작업과 손실은 각 계층의 모델이 적합하도록 보장하며, 이는 모델이 간결한 표현을 생성하도록 보장한다.
\end{remark}






이제 제안된 네트워크 \textsc{crate}의 경험적 성능을 ImageNet-1K에서의 top-1 분류 정확도와 여러 널리 사용되는 다운스트림 데이터셋에서의 전이 학습 성능을 측정하여 제시한다.
결과는 \Cref{tab:crate_comparison_with_sota}에 요약되어 있다. 전이 학습 방법론은 사전 훈련된 네트워크에서 초기화하여 교차 엔트로피 손실을 사용하여 미세 조정하는 것이다.
설계된 화이트박스 트랜스포머 아키텍처는 어텐션 블록(\texttt{MSSA})과 비선형성 블록(\texttt{ISTA}) 모두에서 매개변수 공유를 활용하므로, \textsc{crate}{-Base} 모델(2,280만)은
ViT-Small(2,205만)~\cite{dosovitskiy2020image}과 유사한 수의 매개변수를 가지며, 동일하게 구성된 ViT-Base(8,654만) 매개변수의 30% 미만이다.
\Cref{tab:crate_comparison_with_sota}에서, 우리는 유사한 수의 모델 매개변수로 제안된 네트워크가 ViT와 유사한 ImageNet-1K 및 전이 학습 성능을 달성하면서 간단하고 원칙적인 설계를 가지고 있음을 발견했다. 더욱이, 동일한 훈련 하이퍼파라미터 집합으로, 우리는 \textsc{crate}에서 유망한 스케일링 동작을 관찰한다---우리는 모델 크기를 확장하여 성능을 지속적으로 향상시킨다. 요약하자면, \textsc{crate}는 우리의 원칙적인 아키텍처를 직접 구현함으로써 실제 대규모 데이터셋에서 유망한 성능을 달성한다. 우리는 최종 응용 프로그램 \Cref{ch:applications}에서 이미지 분류에 대한 실험 결과의 구현 및 분석에 대한 자세한 내용을 제공할 것이다.

\begin{figure}[t]
    \begin{subfigure}[t]{0.48\textwidth}
        \centering
        \includegraphics[width=\textwidth]{\toplevelprefix/chapters/chapter4/figs/forward.png}
        \caption{\bf 순방향 전파}
    \end{subfigure}
    \hfill
    \begin{subfigure}[t]{0.48\textwidth}
        \centering
        \includegraphics[width=\textwidth]{\toplevelprefix/chapters/chapter4/figs/backward.png}
        \caption{\bf 역전파}
    \end{subfigure}
    \caption{\small {\bf 심층 네트워크에서 순방향 전파와 역전파의 역할}. (a) 고정된 부분 공간과 사전 $\{(\bm U_{[K]}^{\ell}, \bm D^{\ell})\}_{\ell=1}^L$이 주어지면, 각 계층은 순방향 전파에서 표현에 대한 압축 및 희소화를 수행한다; (b) 역전파는 훈련 데이터로부터 부분 공간과 사전 $\{(\bm U_{[K]}^{\ell}, \bm D^{\ell})\}_{\ell=1}^L$을 학습한다.}
    \label{fig:forward-backward}
\end{figure}


\begin{table*}[t!]
\centering
\caption{\small ImageNet-1K에서 사전 훈련되었을 때 다양한 모델 스케일에서 여러 데이터셋에 대한 \textsc{crate}의 Top-1 분류 정확도. ImageNet-1K/ImageNet-1K ReaL의 경우, top-1 정확도를 직접 평가한다. 다른 데이터셋의 경우, ImageNet에서 사전 훈련된 모델을 초기화로 사용하고 미세 조정을 통해 전이 학습 성능을 평가한다.}
\label{tab:crate_comparison_with_sota}
\small
    \setlength{\tabcolsep}{13.6pt}
\resizebox{0.98\textwidth}{!}{\begin{tabular}{@{}lcccc|cc@{}}
\toprule
\textbf{Model} & \textsc{crate}{-T}  &  \textsc{crate}{-S} & \textsc{crate}{-B} & \textsc{crate}{-L} & { \color{gray} ViT-T} &  { \color{gray}ViT-S } \\ 
\midrule
\midrule
 \# 매개변수 & 6.09M & 13.12M & 22.80M & 77.64M & { \color{gray} 5.72M} & { \color{gray} 22.05M} \\
\midrule
 ImageNet-1K & 66.7 & 69.2 & 70.8 & 71.3 & { \color{gray} 71.5} & { \color{gray} 72.4} \\
 ImageNet-1K ReaL & 74.0 & 76.0 & 76.5 & 77.4 & { \color{gray} 78.3 } & { \color{gray} 78.4} \\
 \midrule
 CIFAR10 & 95.5 & 96.0 & 96.8 & 97.2 & { \color{gray} 96.6} & { \color{gray} 97.2} \\
 CIFAR100 & 78.9 & 81.0 & 82.7 & 83.6 & { \color{gray} 81.8} & { \color{gray} 83.2}\\
 Oxford Flowers-102 & 84.6 & 87.1 & 88.7 & 88.3 & { \color{gray} 85.1} & { \color{gray} 88.5}\\
 Oxford-IIIT-Pets & 81.4 & 84.9 & 85.3 & 87.4 & { \color{gray} 88.5} & { \color{gray} 88.6} \\
 \bottomrule
\end{tabular}}
\end{table*}

\section{설계에 의한 심층 아키텍처 변형} \label{sec:chap4-derive-white-box-transformer-variants}

지금까지 우리는 (인기 있는) 심층 네트워크의 역할이 학습된 표현의 코딩 레이트를 최소화(또는 정보 이득을 최대화)하기 위한 특정 최적화 알고리즘을 실현하는 것이라는 설득력 있는 증거를 제공했기를 바란다. 그러나 최적화 방법에 익숙한 독자들은 위의 아키텍처(ReduNet 또는 CRATE)가 다소 기본적인 최적화 기법에 해당한다는 것을 알아차렸을 것이다. 효율성이나 효과성 면에서 개선의 여지가 많을 수 있다. 더욱이, 우리가 제안한 심층 네트워크 해석을 위한 이론적 프레임워크가 옳다고 믿는다면, 기존 아키텍처를 설명하는 데 도움이 될 뿐만 아니라 더 효율적이고 효과적인 아키텍처를 개발하는 데 지침이 되어야 한다. 이 섹션에서는 이것이 사실일 수 있음을 보여준다: 결과적인 새로운 아키텍처는 완전히 해석 가능할 뿐만 아니라 보장된 정확성과 향상된 효율성을 갖는다. 





\subsection{어텐션 전용 트랜스포머 아키텍처} \label{sub:aot}

이 하위 섹션에서는 MSSA 연산자를 기반으로 하는 해석 가능한 계층으로 구성된 최소한의 트랜스포머 아키텍처를 제안한다. 필요한 구성 요소만으로 완전히 해석 가능한 트랜스포머 아키텍처를 도출하기 위해,
우리는 표현 학습의 목표가 잡음이 있는 초기 토큰 표현 집합을 저차원 부분 공간의 혼합으로 압축하는 것이라고 주장한다. %
여기서, 우리는 초기 토큰 표현 $\bm Z^{(1)}$이 다음과 같이 잡음으로 교란된 저차원 가우시안의 혼합에서 샘플링된다고 가정한다:
\begin{definition}\label{def:MoG}
$C_1,\dots,C_K$를 인덱스 집합 $[N]$의 분할이라 하고, 각 $k \in [K]$에 대해 $\bm U_k \in \mathcal{O}^{d \times p_k}$가 $k$-번째 부분 공간의 직교 정규 기저를 나타낸다고 하자. 우리는 토큰 표현 $\{\bm z_i \}_{i=1}^N \subseteq \R^d$이 각 $k \in [K]$에 대해 다음과 같은 잡음이 있는 저차원 가우시안 분포의 혼합에서 샘플링되었다고 말한다.
\begin{align}\label{eq:MoG}
    \bm z_i = \underbrace{\bm U_{k} \bm a_i}_{\bf 신호} + \underbrace{\sum_{j \neq k}^K \bm U_j \bm e_{i,j}}_{\bf 잡음},\ \forall i \in C_k, 
\end{align}
여기서 모든 $i \in C_k$와 $k \in [K]$에 대해 $\bm{a}_i \overset{i.i.d.}{\sim} \mathcal{N}(\bm{0},\bm{I}_{p_k})$이고 $\bm{e}_{i,j} \overset{i.i.d.}{\sim} \mathcal{N}(\bm{0},\delta^2\bm{I}_{p_j})$이며, $\{\bm{a}_i\}$와 $\{\bm{e}_{i,j}\}$는 각각 상호 독립이고, $\{\bm{a}_i\}$는 $\{\bm{e}_{i,j}\}$와 독립이다.         
\end{definition}
이 모델은 실제 사전 훈련된 LLM에서 토큰 표현을 근사하기 위한 이상적인 프레임워크 역할을 한다. 이 모델은 토큰 표현이 잡음이 있는 여러 저차원 가우시안 분포의 혼합에서 샘플링된다고 가정한다. 이 모델 하에서, 표현 학습의 목표는 잡음이 있는 초기 토큰 표현 집합을 해당 부분 공간으로 압축하는 것이다. 또한, 이 모델은 사전 훈련된 대규모 언어 모델의 토큰 표현 구조에 대한 두 가지 잘 확립된 가설인 "선형 표현 가설" \citep{jiang2024origins,park2023linear}과 "중첩 가설" \citep{elhage2022toy,yun2021transformer}과 잘 일치한다.

\begin{remark}
    선형 표현 가설은 LLM의 토큰 표현이 의미론적 특징을 인코딩하는 저차원 선형 부분 공간에 놓여 있다고 가정한다. 유사하게, 중첩 가설은 이러한 표현이 이러한 특징 벡터의 희소 선형 조합으로 근사적으로 표현될 수 있다고 제안한다. \Cref{def:MoG}에서, 부분 공간의 각 기저 $\bm U_k$는 의미론적 특징 집합으로 해석될 수 있으며, 각 특징은 토큰 의미의 특정 측면에 해당한다. 그런 다음 토큰 표현은 이러한 부분 공간 기저의 희소 선형 조합으로 근사적으로 표현되어, 관련 없는 차원을 무시하면서 토큰의 본질적인 의미론적 구성 요소를 포착한다. 
\end{remark}

\paragraph{토큰 표현을 위한 잡음 제거 연산자.} 이제, MSSA 연산자(\eqref{eq:Multi-Head-SSA} 참조)가 위 모델에서 생성된 토큰 표현의 잡음을 점진적으로 제거할 수 있음을 보여준다. 구체적으로, 우리는 각 $\ell =1 ,\dots,L$에 대해 다음을 고려한다. 
\begin{align}\label{eq:MSSA}
    \bm Z^{(\ell+1)} =  \bm Z^{(\ell)} + \eta \sum_{k=1}^K \bm U_k\bm U_k^T \bm Z^{(\ell)} \varphi \left(\bm Z^{(\ell)^T}\bm U_k\bm U_k^T\bm Z^{(\ell)} \right),
\end{align}
여기서 $\{\bm U_k\}_{k=1}^K$는 \Cref{def:MoG}에 정의되어 있고, $\eta > 0$는 스텝 크기이며, $\varphi$는 소프트맥스, ReLU 또는 다른 함수와 같은 원소별 연산자이다. 우리의 개발을 단순화하기 위해, \Cref{def:MoG}의 부분 공간들이 서로 직교한다고 가정한다. 즉, 모든 $k \neq j$에 대해 $\bm U_k^T\bm U_j = \bm 0$이다. 고차원 공간에서는 무작위 저차원 부분 공간이 높은 확률로 서로 비일관적이므로, 즉 $\bm U_k^T\bm U_j \approx \bm 0$이므로 \citep{Wright-Ma-2021}, 이 가정은 제한적이지 않다는 점에 유의하라.\footnote{약간 더 정교한 분석을 통해 우리의 결과를 비직교 부분 공간으로 간단히 일반화할 수 있다.}   

\begin{figure}[t]
    \begin{subfigure}[t]{0.45\textwidth}
        \centering
        \includegraphics[width=\textwidth]{\toplevelprefix/chapters/chapter4/figs/SNR1.png}
        \caption{잡음 수준 $\delta = 0.2$}
    \end{subfigure}
    \hfill
    \begin{subfigure}[t]{0.45\textwidth}
        \centering
        \includegraphics[width=\textwidth]{\toplevelprefix/chapters/chapter4/figs/SNR2.png}
        \caption{잡음 수준 $\delta = 0.5$}
    \end{subfigure}
    \caption{{\bf 어텐션 전용 트랜스포머의 잡음 제거 성능.} 여기서, 우리는 \Cref{def:MoG}의 저차원 가우시안 혼합에서 초기 토큰 표현을 샘플링한다. 그런 다음, \eqref{eq:MSSA}를 적용하여 토큰 표현을 업데이트하고 각 계층의 SNR을 보고한다.}  \label{fig:MSSA} 
\end{figure} 


이제, $\bm Z_k^{(\ell)}$의 열을 $\ell$-번째 계층에서 $k$-번째 부분 공간으로부터의 토큰 표현이라고 하자. 잡음 제거 능력을 정량화하기 위해, 우리는 $\ell$-번째 계층에서 토큰 표현의 각 블록에 대한 신호 대 잡음비(SNR)를 다음과 같이 정의한다:
\begin{align}\label{def:SNR}
\mathrm{SNR}(\bm Z_k^{(\ell)}) \doteq  \frac{\|\bm U_k\bm U_k^T\bm Z_k^{(\ell)} \|_F}{\|(\bm I - \bm U_k\bm U_k^T)\bm Z_k^{(\ell)} \|_F},\quad \forall k \in [K].
\end{align}
분석을 단순화하기 위해, $p=p_1=\dots=p_K$, $N_1=\dots=N_K=N/K$라고 가정하고,
\begin{align}\label{eq:orth}
\begin{bmatrix}
\bm U_1 & \dots & \bm U_K
\end{bmatrix} \in \mathcal{O}^{d\times Kp}. 
\end{align}  

위의 설정으로, 이제 MSSA 연산자의 잡음 제거 성능을 특성화한다.

\begin{figure*}[t]
\begin{center}
        \includegraphics[width=0.75\textwidth]{\toplevelprefix/chapters/chapter4/figs/MSSA.png}
    \caption{\textbf{어텐션 전용 트랜스포머 아키텍처의 세부 사항.} 각 계층은 MSSA 연산자와 스킵 연결로 구성된다. 또한, 언어 작업에 대해서는 LayerNorm이 포함된다. 실제로는 훈련 샘플을 사용하여 모델 매개변수를 훈련하기 위해 역전파가 적용된다. %
    }\label{fig:transformer}
\end{center} 
\end{figure*}

\begin{theorem}\label{thm:1}
$\bm Z^{(1)}$이 \Cref{def:MoG}에 정의되고 \eqref{eq:MSSA}의 $\varphi(\cdot)$가
$\varphi(\bm x) = h\left(\sigma(\bm x)\right)$라고 하자.
여기서 $\sigma:\R^N \to \R^N$는 소프트맥스 함수이고 $h:\R^N \to \R^N$는 각 $i \in [N]$에 대해 $h(x) = \tau \mathbb{I}\left\{x > \tau\right\}$인 원소별 임계값 함수이다. $p \gtrsim \log N$, $\delta \lesssim \sqrt{\log N}/\sqrt{p}$이고,
\begin{align*}
\tau \in \left( \frac{1}{2},  \frac{1}{1+N\exp(-9p/32)} \right].
\end{align*}
충분히 큰 $N$에 대해, 각 $\ell \in [L]$에 대해 확률적으로 최소 $1-KLN^{-\Omega(1)}$로 다음이 성립한다.
    \begin{align}\label{eq:SNR}
        \mathrm{SNR}(\bm Z_k^{(\ell+1)}) = (1+\eta\tau) \mathrm{SNR}(\bm Z_k^{(\ell)}),\ \forall k \in [K]. 
    \end{align}
\end{theorem}
이 정리는 초기 토큰 표현이 잡음 수준 $O(\sqrt{\log N}/\sqrt{p})$인 저차원 가우시안 분포의 혼합에서 샘플링될 때, 제안된 트랜스포머의 각 계층이 선형 속도로 토큰 표현의 잡음을 제거함을 보여준다. 이는 MSSA 연산자가 계층에 걸쳐 잡음을 줄이는 데 효율적임을 나타낸다. 특히, 우리의 이론적 결과는 \Cref{fig:MSSA}의 실험적 관찰에 의해 잘 뒷받침된다. 이 정리는 (\ref{eq:MSSA})를 펼침으로써 파생된 트랜스포머 아키텍처의 실제적인 잡음 제거 능력에 대한 이론적 기반을 제공한다.

\begin{remark}
    이 모델 하에서, 표현 학습의 목표는 잡음이 있는 초기 토큰 표현 집합을 해당 부분 공간으로 압축하는 것이다. 그러나 실제 응용에서는 토큰 표현이 더 복잡한 구조를 나타내므로, 표현 학습의 목표는 토큰 집합을 압축하여 간결하고 구조화된 표현을 찾는 것임을 지적해야 한다.
\end{remark}


\paragraph{어텐션 전용 트랜스포머.} 이제, 어텐션 전용 트랜스포머 아키텍처를 공식적으로 제안한다. 구체적으로, 반복적인 최적화 단계 (\ref{eq:MSSA})를 심층 네트워크의 계층으로 펼침으로써, \Cref{fig:transformer}에서 트랜스포머 아키텍처를 구성한다. 제안된 아키텍처의 각 계층은 MSSA 연산자와 스킵 연결만으로 구성된다. 언어 작업의 경우, 성능 향상을 위해 MSSA 연산자 앞에 LayerNorm을 추가로 통합한다. 완전한 아키텍처는 위치 인코딩, 토큰 임베딩 및 최종 작업별 헤드와 같은 필수적인 작업별 전처리 및 후처리 단계와 함께 이러한 계층을 쌓아서 구축된다.




일반적으로, 표준 디코더 전용 트랜스포머 아키텍처는 다음과 같은 주요 구성 요소로 구성된다 \cite{vaswani2017attention}: (1) 위치 인코딩, (2) 다중 헤드 QKV 자기-어텐션 메커니즘, (3) 피드포워드 MLP 네트워크, (4) 계층 정규화, (5) 잔차 연결. 이와 대조적으로, 제안된 트랜스포머 아키텍처는 몇 가지 주요 단순화를 통합하여 간소화된 설계를 채택한다. 구체적으로, 공유-QKV 부분 공간 자기-어텐션 메커니즘을 사용하고, MLP 계층을 제외하며, LayerNorm의 빈도를 줄인다.










\subsection{선형 시간 어텐션: 토큰 통계 트랜스포머}\label{sub:tost} 

이 하위 섹션에서는 코딩 레이트 감소 목표를 기반으로 토큰 수에 선형적으로 확장되는 계산 복잡도를 가진 새로운 트랜스포머 어텐션 연산자를 제안한다. 구체적으로, 우리는 MCR$^2$ 목표의 새로운 변형 형식을 유도하고 이 변형 목표의 펼쳐진 경사 하강법에서 비롯된 아키텍처가 토큰 통계 자기 어텐션(\texttt{TSSA})이라는 새로운 어텐션 모듈로 이어진다는 것을 보여준다. \texttt{TSSA}는 {\em 선형 계산 및 메모리 복잡도}를 가지며 토큰 간의 쌍별 유사성을 계산하는 일반적인 어텐션 아키텍처와 근본적으로 다르다. \eqref{eq:MCR pi}에서 $\bm \Pi = [\bm \pi_1, \ldots, \bm \pi_K] \in \R^{N \times K}$는 확률적 "그룹 할당" 행렬을 나타내며(즉, $\bm \Pi \bm 1 = \bm 1$이고 모든 $(i,k) \in [N] \times [K]$에 대해 $\Pi_{ik} \geq 0$), 여기서 $\Pi_{ik}$는 $i$-번째 토큰을 $k$-번째 그룹에 할당할 확률을 나타낸다.

 

\paragraph{코딩 레이트에 대한 새로운 변형 형식.} 시작하기 위해, 우리는 행렬의 스펙트럼의 오목 함수를 기반으로 한 MCR$^2$ 유사 목표의 일반적인 형태를 고려한다. 즉, 주어진 PSD 행렬 $\bm M \in \PSD(d)$와 임의의 스칼라 $c \geq 0$에 대해 $\log \det (\I + c \bm M) = \sum_{i=1}^{d} \log( 1 + c \lambda_i(\bm M))$이며, 여기서 $\lambda_i (\bm M)$는 $\bm M$의 $i$-번째로 큰 고유값이다. 또한 $\log(1 + c \sigma)$는 $\sigma$의 오목 비감소 함수임에 유의하라. 따라서, 우리는 $f$가 오목하고 비감소적인 $F(\bm M) = \sum_{i=1}^{d} f(\lambda_i(\bm M))$ 형태의 PSD 행렬의 일반적인 스펙트럼 함수를 기반으로 한 MCR$^2$의 더 일반적인 형태로 결과를 설명한다. 특히, 위 논의에서 어텐션 메커니즘이 MCR$^2$의 압축 구성 요소를 펼침으로써 발생한다는 것을 상기하고, 더 일반적인 MCR$^2$ 스타일 압축 함수를 고려한다:
\vspace{-2mm}
\begin{equation}
    R_{c,f}(\bm Z, \bm \Pi) \doteq \frac{1}{2}\sum_{k=1}^K \frac{N_k}{N} F\left( \frac{1}{N_k}\Z \mathrm{Diag}(\bm \pi_k) \Z^{\top} \right).
    \label{eq:mcr_c_gen}
\end{equation}



위의 목표에 대해, 우리는 이제 다음 결과를 주목한다:
\begin{theorem}
\label{thm:var_concave}
    \(f \colon [0, \infty) \to \R\)가 비감소, 오목하고 \(f(0) = 0\)을 만족한다고 하자. 그리고 \(F \colon \PSD(d) \to \R\)가 \(F(\bm M) = \sum_{i = 1}^{d}f(\lambda_{i}(\bm M))\) 형태를 갖는다고 하자. 그러면 각 \(\bm M \in \PSD(d)\)와 \(\bm Q \in \O(d)\)에 대해 다음이 성립한다.
    \begin{equation}
        \label{eq:U_bound}
        F(\bm M) \leq  \sum_{i=1}^{d} f\left( (\bm Q^{\top} \bm M \bm Q)_{ii} \right).
    \end{equation}
    또한, \eqref{eq:U_bound}의 부등식은 $\bm M$을 대각화하는 임의의 $\bm Q$에 대해 등호로 달성되며, $f$가 엄격하게 오목하면 \eqref{eq:U_bound}의 부등식은 $\bm Q$가 $\bm M$을 대각화하는 경우에만 등호로 달성된다.
\end{theorem}

위 결과를 사용하여 \eqref{eq:mcr_c_gen}을 다음과 같은 형태의 동등한 변형 목표로 대체할 수 있다.
\vspace{-2mm}
\begin{equation}
    \label{eq:var_comp}
    R^{\rm var}_{c,f} (\Z,\bm \Pi \mid \bm U_{[K]}) \doteq \frac{1}{2}\sum_{k=1}^K \frac{N_k}{N} \sum_{i=1}^{d} f\left(\frac{1}{N_k} (\bm U_{k}^{\top} \Z \mathrm{Diag}(\bm \pi_k) \bm Z^{\top} \bm U_{k})_{ii} \right),
\end{equation}
여기서 동등성은 \Cref{thm:var_concave}에서 설명된 $\{ \bm U_k \in \O(d)\}_{k=1}^K$ 행렬의 최적 선택(즉, 각 $\Z \mathrm{Diag}(\bm \pi_k) \Z^\top$을 대각화하는 직교 행렬)에 대해 $ R^{\rm var}_{c,f} (\Z,\bm \Pi \mid \bm U_{[K]}) = R_{c,f} (\Z,\bm \Pi)$로 엄격한 경계를 달성한다는 의미이다. 일반적으로, 이 경계를 달성하려면 각 샘플링된 $\Z$ 인스턴스에 대해 각 $\Z\mathrm{Diag}(\bm \pi_{k})\Z^{\top}$을 대각화하는 새로운 최적의 $\bm U_{k}$ 매개변수 집합을 선택해야 하며, 이는 네트워크 아키텍처에 대해 명백히 비실용적이다.
대신, 대안적인 관점으로서, 데이터($\Z$)를 고정하고 엄격한 변형 경계를 달성하기 위해 $\bm U_k$ 매개변수를 최적화하는 대신, 위에서 설명한 알고리즘적 펼침 설계 원칙을 적용하고 $\Z$를 섭동하여 $R_{c, f}^{\rm var}(\cdot \mid \bm U_{[K]})$를 점진적으로 최소화하는 연산자를 설계할 수 있다. 이 점을 명확히 하기 위해, 각 변형 경계는 $\Z \mathrm{Diag}(\bm \pi_k) \Z^\top$의 고유 공간이 $\bm U_k$의 열과 정렬될 때 엄격해지므로, $\Z$의 적절한 열(즉, $\bm \pi_k$의 큰 항목에 해당하는 열)을 회전하여 $\bm U_k$와 정렬함으로써 엄격한 변형 경계에 접근할 수 있다. 즉, 각 $\Z$ 인스턴스에 대해 데이터를 정렬하기 위해 $\bm U_k$를 회전하는 대신, 각 $\Z$의 토큰 특징을 회전하여 $\bm U_k$와 정렬할 수 있다.

이 접근 방식에 따라, 우리는 $R_{c,f}^{\rm var}$에 대한 $\Z$에 대한 경사 하강 단계를 계산한다.
이 계산을 시작하기 위해, 먼저 $\bm \pi \in \Re^N$을 임의의 원소별 음수가 아닌 벡터라고 하자. 그러면 다음을 얻는다.
\begin{equation}\label{eq1:grad Rcf}
\nabla_{\Z} \ \frac{1}{2} \sum_{i = 1}^{d} f(  (\Z \mathrm{Diag}(\bm \pi) \Z^{\top} )_{ii}) %
= \; \mathrm{Diag}(\nabla f[ \Z^{\hada 2} \bm \pi] ) \Z \mathrm{Diag}(\bm \pi),
\end{equation}
여기서 $\nabla f$는 $f$의 그래디언트이고, (상기) $\nabla f[\cdot]$는 브래킷 안의 벡터의 각 원소에 $\nabla f$를 적용한다. 특히, $ f(x) = \log(1 + (d/\epsilon^{2}) x)$의 경우, $\nabla f(x) = (d / \epsilon^{2}) (1+ (d / \epsilon^{2}) x)^{-1}$는 단순히 비선형 활성화이다. 또한, (상기) $N_{k} = \langle \bm \pi_{k}, \bm 1\rangle$이다. 따라서, $R^{\rm var}_{c,f}$의 $\Z$에 대한 그래디언트는 다음과 같다:
\begin{align}\label{eq2:grad Rcf}
    \nabla_{\Z} R^{\rm var}_{c,f}(\Z,\bm \Pi \mid \bm U_{[K]}) =  \frac{1}{n}
    \sum_{k=1}^K \bm U_{k} \underbrace{\mathrm{Diag} \left( \nabla f\left[ (\bm
    U_{k}^{\top} \Z)^{\hada 2}  \frac{\bm \pi_k}{\langle \bm \pi_{k}, \bm 1\rangle} \right] \right)}_{\doteq \bm D(\Z, \bm \pi_{k} \mid \bm U_{k})} \bm U_{k}^{\top} \Z \mathrm{Diag}(\bm \pi_k).
\end{align}
($1/N$ 상수는 합의 각 항에 있는 $(N_{k}/N)\cdot (1/N_{k}) = 1/N$ 상수에서 비롯된다는 점에 유의하라.) 이제 $j$-번째 토큰 $\z_j$에 대한 경사 단계를 고려하면
, 우리는 제안된 증분 압축 연산자, 즉 \textit{자기 어텐션} + 잔차 연산자의 대리자에 도달한다:
\vspace{-2mm}
\begin{equation}\label{eq:layer-operation-orig}
    \z_j^+ = \z_{j} - \tau \nabla_{\z_{j}}R_{c, f}^{\rm var}(\Z, \bm \Pi \mid \bm U_{[K]}) = \z_j - \frac{\tau}{N} \sum_{k=1}^K \Pi_{jk} \bm U_{k} \bm D(\Z, \bm \pi_{k} \mid \bm U_{k}) \bm U_{k}^{\top} \z_j
\end{equation}
각각의 \(j \in [n]\)에 대해, 여기서 $\tau > 0$는 증분 최적화를 위한 스텝 크기 매개변수이다. 그런 다음, \Cref{fig:vcrate-architecture}에서 TOST의 계층을 구성할 수 있다.

\begin{figure}[t]
    \centering \includegraphics[width=1\textwidth,trim={0 5.0cm 0 0},clip]{\toplevelprefix/chapters/chapter4/figs/V-CRATE.pdf}
    \vspace{-1mm}
    \caption{\small \textbf{제안된 토큰 통계 트랜스포머(ToST)의 한 계층 $\ell$.} 특히, ToST의 자기 어텐션은 투영된 토큰의 각 행에 \textit{오직 스칼라}만을 곱하여 토큰 $\Z^{\ell}$을 효율적으로 $\Z^{\ell+1}$로 변환한다. 이는 어텐션의 복잡도를 감소시킨다: 각 헤드의 투영된 토큰 차원이 $p$이고 토큰 수가 $n$일 때, $O(p)$ 공간 및 $O(pn)$ 시간 복잡도를 갖는다.
    }
    \label{fig:vcrate-architecture}
\end{figure}

\paragraph{모델 해석.} 제안된 어텐션 연산자 \eqref{eq:layer-operation-orig}가 주어지면, 먼저 $\bm\Pi$의 행이 음수가 아니고 합이 1임을 상기하자.
, 그래서 우리 연산자는 $K$개의 "어텐션 헤드" 같은 연산자들의 가중 평균을 취한 다음 잔차 연결을 추가한다. \(\sum_{k = 1}^{K}\Pi_{jk} = 1\)을 사용하여, \eqref{eq:layer-operation-orig}를 다음과 같이 다시 쓸 수 있다: %
\vspace{-2mm}
\begin{equation}
\label{eq:z_up}
    \z_j^+ = \sum_{k=1}^K \Pi_{jk} \Big[ \z_j \underbrace{- \frac{\tau}{n} \bm U_{k} \bm D(\Z, \bm \pi_{k} \mid \bm U_{k}) \bm U_{k}^{\top}}_{\text{하나의 어텐션 헤드의 작용}} \z_j \Big].
\end{equation}
즉, 각 어텐션 헤드를 먼저 토큰 특징을
$\bm U_k^\top$를 곱하여 기저 $\bm U_{k}$에 투영하고, 대각 행렬
$\bm D(\Z, \bm \pi_{k} \mid \bm U_{k})$(줄여서 \(\bm
D_{k}\))를 곱하고, $\bm
U_{k}$를 곱하여 표준 기저로 다시 투영한 다음, 잔차
연결을 통해 원래 토큰 특징에서 이것을 빼는 것으로 볼 수 있다. 우리 어텐션 계층의 핵심 측면은 $\bm
D_{k}$의 계산이다. 즉, \(\Pi_{jk} \geq 0\)이므로 $\bm \pi_k / \langle \bm \pi_{k}, \bm
1\rangle \in \Re^N$는 토큰이
$k^\text{th}$ 그룹에 속하는 확률 분포를 형성한다. 결과적으로, $(\bm U^\top_k \Z)^{\hada 2} \bm \pi_k / \langle \bm \pi_{k}, \bm 1\rangle$는 $\bm \pi_k /  \langle \bm \pi_{k}, \bm 1\rangle$에 의해 주어진 분포 하에서 $\bm U_k^\top \Z$의 2차 모멘트를 추정한다. 또한, $f$가 오목 비감소 함수이므로, $\nabla f(x)$는 $x$가 증가함에 따라 단조롭게 $0$으로 감소하므로, $\bm D_{k}$의 항목들($\nabla f(x)$ 형태를 가짐)은 $x=0$에서 최댓값을 달성하고 %
$x$가 증가함에 따라 단조롭게 $0$으로 감소한다.

이로부터, 우리는 어텐션 헤드 + 잔차
연산자 $[\I - (\tau/n)\bm U_{k}\bm D_{k}\bm U_{k}^{\top}]$의 핵심 해석에 도달한다. 즉, 이
연산자는 근사적인 저차원 데이터 종속 투영을 수행하며, 여기서 투영 $\bm
U_k^\top \Z$ 후 "파워"가 큰 방향(즉, 2차 모멘트 $(\bm
U_{k}^{\top}\Z)^{\hada 2}\bm \pi_k / \langle \bm \pi_{k}, \bm 1\rangle$가 큰 방향)은 보존되는 반면, 그렇지 않은 방향은 억제된다. 이를 보기 위해, $\bm D_k$의 항목들은 2차 모멘트가 증가함에 따라 단조롭게 0으로 감소하므로, 큰 2차 모멘트를 가진 방향에 대해 어텐션 + 잔차 연산자는 대체로 항등 연산자처럼 작동한다는 것을 상기하자. 반대로, 작은 2차 모멘트를 가진 방향에 대해, 우리 연산자는 해당 방향을 따라 토큰의 투영을 빼서, 결과적으로 해당 방향이 억제된다. 표준 자기-어텐션 연산자와 비교할 때, 우리 방법은 명확하게 토큰 간의 쌍별 유사성을 계산하지 않는다. 오히려, $\Z$의 토큰 간의 상호 작용은 $\bm D_{k}$를 구성하는 데 사용되는 2차 모멘트 통계에 대한 기여를 통해서만 연산자에 영향을 미친다. 그럼에도 불구하고, 표준 자기-어텐션 연산자와 유사하게, 우리 방법은 여전히 데이터 종속 저차원 구조로의 압축 형태를 수행하는 것으로 명확하게 해석될 수 있다. 즉, 억제되는 특정 방향이 그룹의 다른 토큰과 강하게 정렬되지 않은 방향인 근사적인 저차원 투영을 수행한다는 의미에서 그렇다.

\paragraph{실용적인 구현 세부 사항.} 제안된 어텐션 연산자를 소개했으므로, 이제 더 실용적인 고려 사항에 대해 논의한다. 첫째, 발표의 이 시점까지, 우리는 $\bm\Pi$ 행렬을 통해 토큰이 어떻게 다양한 어텐션 헤드로 "그룹화"되는지에 대한 논의를 피했지만, 우리 방법을 구현하려면 $\bm\Pi$를 구성하는 수단이 명확히 필요하다. 또한, \Cref{thm:var_concave}의 변형 형식은 $\bm U$ 행렬이 정사각형이고 직교해야 하지만, 효율성을 위해 더 작은 행렬(즉, $\bm U$의 열 수를 줄임)을 사용하고 직교성 제약 조건을 제거하는 것이 이상적일 것이다.

실제로, 우리는 직교성 제약 조건을 강제하지 않는다. $\bm U$ 행렬의 열 수를 줄이기 위해, CRATE \citep{yu2023white}와 유사하게, 그룹 \(k\) 내의 특징 $\bm Z$이 저차원 부분 공간(차원 \(p\)) 주위에 (근사적으로) 군집되어 있다고 가정하면, 그룹 내 \(k\) 공분산 \(\Z\mathrm{Diag}(\bm \pi_{k})\Z^{\top}\)는 저차원이며, \cite{yu2020learning}은 $\Z$의 최적 기하학이 각 그룹이 다른 그룹과 직교하는 저차원 부분 공간이 되는 것임을 보여준다. 따라서 우리는 이 공분산의 이미지, 즉 그룹 \(k\)의 데이터의 선형 생성에 대한 저차원 직교 정규 기저를 명시적으로 찾을 수 있다. 차원이 $p \leq d$이면, 기저는 $d\times p$ 직교 행렬 $\bm U_k \in \O(d, p)$로 표현될 수 있다. 이 경우, 이러한 저차원 직교 기저 행렬을 사용하여 \(R_{c,f}\)를 더 효율적으로 상한할 수 있다. 이를 보이기 위해, \Cref{thm:var_concave}의 더 일반적인 버전을 사용하여 다음 보조 정리를 도출한다.
\begin{corollary}\label{cor:var_concave_logdet}
    \(f \colon [0, \infty) \to \R\)가 비감소, 오목하고 \(f(0) = 0\)을 만족한다고 하자. 그리고 \(F \colon \PSD(p) \to \R\)가 \(F(\bm M) = \sum_{i = 1}^{p}f(\lambda_{i}(\bm M))\) 형태를 갖는다고 하자. \(\Z\), \(\bm \Pi\)가 고정되었다고 하자. 그러면, 모든 \(k \in [K]\)에 대해 \(\mathrm{image}(\Z\diag(\bm \pi_{k})\Z^{\top}) \subset \mathrm{image}(\bm U_{k})\)인 모든 \(\bm U_{1}, \dots, \bm U_{K} \in \O(d, p)\)에 대해 다음이 성립한다.
    \begin{equation}
        R_{c, f}(\Z, \bm\Pi) \leq R_{c, f}^{\rm var}(\Z, \bm\Pi \mid \bm U_{[K]}), 
    \end{equation}
     여기서 \(R_{c, f}^{\rm var}\)는 \eqref{eq:var_comp}에 공식적으로 정의되어 있다. 등호는 각 \(k \in [K]\)에 대해 \(\bm U_{k}\)가 \(\Z\diag(\bm \pi_{k})\Z^{\top}\)를 대각화하는 경우에 성립하며, \(f\)가 강하게 오목하면 이 등호 조건은 "필요충분조건"이 된다.
\end{corollary}

우리의 어텐션 연산자를 정의하는 마지막 단계는 그룹 소속 $\bm\Pi$를 추정하는 것이다. 이를 위해 우리는 각 특징 \(\z_{j}\)가 지지하는 부분 공간에서 어떻게 벗어나는지에 대한 간단한 모델을 가정하고 최적의 부분 공간 할당을 찾는다. \cite{yu2023white}는 각 \(\z_{j}\)를 독립적으로 저차원 가우시안 혼합 모델에 속하는 것으로 모델링하면, 각 가우시안이 동일한 스펙트럼을 가진 공분산 행렬을 갖고 직교 정규 기저 \(\bm U_{k}\)를 가진 부분 공간에 지지되며, 공분산 \(\eta \I\)를 가진 독립적인 가우시안 잡음이 더해진다고 가정하면, 각 토큰 \(\z_{j}\)가 각 부분 공간에 속할 사후 확률이 다음과 같이 할당 행렬 \(\bm \Pi = \bm \Pi(\bm Z \mid \bm U_{[K]})\)에 의해 주어진다는 것을 보여준다:
\begin{align}\label{eq:pi}
    \bm \Pi = \begin{bmatrix} \boldsymbol{\nu}(\z_1 \mid \bm U_{[K]})^\top \\ \vdots \\ \boldsymbol{\nu}(\z_n \mid \bm U_{[K]})^\top \end{bmatrix}, \quad
\text{where} \quad 
\boldsymbol{\nu}(\z_j \mid \bm U_{[K]}) \doteq \operatorname{softmax}\left( \frac{1}{2\eta} \begin{bmatrix} \|  \bm U_{1}^{\top} \z_j \|_2^2 \\ \vdots \\ \| \bm U_{K}^{\top} \z_j \|_2^2 \end{bmatrix} \right), \quad \forall j \in [n],
\end{align}
여기서 $\eta$는 학습 가능한 온도 매개변수가 된다. 따라서, 입력 특징 \(\Z\)가 주어지면, \eqref{eq:pi}를 사용하여 \(\bm\Pi\)를 추정하고 어텐션 연산자를 계산한다. \eqref{eq:pi}의 $\bm\Pi$ 구성과
\eqref{eq:layer-operation-orig}를 결합하면, {\em 토큰 통계 자기-어텐션} 연산자를 얻는다:
\begin{equation}
    \label{eq:tssa}
   \texttt{TSSA}(\Z \mid \bm U_{[K]}) \doteq -\frac{\tau}{n}\sum_{k = 1}^{K}\bm U_{k}\bm D(\Z, \bm\pi_{k} \mid \bm U_{k})\bm U_{k}^{\top}\Z\diag(\bm \pi_{k}),
\end{equation}
여기서 \(\bm\pi_{k}\)는 \eqref{eq:pi}에 정의된 \(\bm\Pi = \bm\Pi(\Z \mid \bm U_{[K]})\)의 열이고 \(\bm D\)는 \eqref{eq2:grad Rcf}에 정의되어 있다. 


\section{요약 및 참고 사항}

이 장에 제시된 자료는 \cite{chan2021redunet, wang2024global, wang2025attention, wu2025token, yu2023white}를 포함한 이 주제에 대한 일련의 최근 연구에 기반을 두고 있다. 이러한 기여는 펼쳐진 최적화를 통해 해석 가능한 심층 네트워크를 구성하기 위한 이론적 발전과 실제적 방법론을 모두 포함한다. 이 장에서 논의된 많은 핵심 결과와 증명은 이러한 기초 연구에서 직접 파생되었거나 영감을 받았다. 


최적화 알고리즘을 펼쳐 신경망을 구성하는 아이디어는 \cite{gregor2010learning}의 독창적인 연구로 거슬러 올라간다. 이 연구에서 저자들은 반복 축소-임계값 알고리즘(ISTA)과 같은 희소 코딩 알고리즘을 펼쳐 다층 퍼셉트론(MLP)을 형성할 수 있음을 보여주었으며, 이는 반복 최적화와 신경망 설계를 효과적으로 연결했다. 특히, \cite{monga2019algorithm}은 이러한 펼쳐진 네트워크가 일반적인 네트워크에 비해 더 해석 가능하고, 매개변수 효율적이며, 효과적임을 입증했다. 이 장에서는 이전에 소개된 (희소) 레이트 감소 목표와 같이 잘 동기 부여된 목표를 최소화하도록 설계된 최적화 알고리즘을 펼침으로써 원칙에 입각한 화이트박스 심층 네트워크 아키텍처를 개발하기 위해 이 관점을 기반으로 한다. 이 접근 방식은 네트워크의 각 계층의 역할을 명확히 할 뿐만 아니라, 경험적 시행착오를 넘어 해석 가능하고 목표 지향적인 설계를 향한 아키텍처 선택에 대한 이론적 근거를 제공한다. 다음에서는 일반적으로 경험적 설계와 휴리스틱 튜닝을 통해 구성되는 기존 DNN과 수학적으로 기반을 둔 ReduNet 아키텍처를 비교한다: 

\begin{center}
\begin{tabular}{| l || c | c |}
\hline
  & 기존 DNN & ReduNets\\ [0.5ex]
  \hline \hline
목표 & 입/출력 맞춤 & 정보 이득\\ [0.5ex]
  \hline
심층 아키텍처 & 시행착오 & 반복 최적화 \\  [0.5ex]
\hline
계층 연산자 & 경험적 & 투영 경사 \\  [0.5ex]
\hline
이동 불변성 & CNN+증강 & 불변 ReduNets \\  [0.5ex]
\hline
초기화 & 무작위/사전 설계 & 순방향 펼침 \\ [0.5ex]
\hline
훈련/미세 조정 & 역전파 & 순방향/역전파\\ [0.5ex]
\hline
해석 가능성 & 블랙박스 & 화이트박스 \\ [0.5ex]
\hline
표현 & 은닉/잠재 & 비일관적 부분 공간 \\ [0.5ex]
\hline
\end{tabular}
\end{center}




\section{연습문제 및 확장}


\begin{exercise}
    $\bm Z = [\bm Z_1,\dots,\bm Z_K] \in \R^{d\times m}$이고 각 $k \in [K]$에 대해 $\bm Z_k \in \R^{d\times m_k}$라고 하자. 어떤 $\alpha > 0$에 대해, 다음을 가정하자.
    \begin{align*}
        R(\bm Z) = \log\det\left(\bm I + \alpha\bm Z\bm Z^T \right).
    \end{align*}
1. 임의의 방향 $\bm D \in \R^{d\times m}$이 주어졌을 때, $\nabla R(\bm Z) = \alpha\bm X^{-1}\bm Z$이고 다음을 보여라.
    \begin{align*}
 \nabla^2 R(\bm Z)[\bm D, \bm D] =\alpha \mathrm{Tr}\left( \bm X^{-1}\bm D\bm D^T\right) - \frac{\alpha^2}{2}\mathrm{Tr}\left(\bm X^{-1}\left( \bm Z\bm D^T+\bm D\bm Z^T\right) \bm X^{-1}\left( \bm Z\bm D^T+\bm D\bm Z^T\right)\right),
    \end{align*}
    여기서 $\bm X \doteq \bm I + \alpha\bm Z\bm Z^T$이다. {\em 힌트:} 다음을 참고하라.
    \begin{align*}
        \nabla^2 R(\bm Z)[\bm D, \bm D] \doteq \left\langle \lim_{t \to 0} \frac{\nabla R(\bm Z+ t\bm D) - \nabla R(\bm Z)}{t}, \bm D \right\rangle. 
    \end{align*}

\noindent
2. 다음을 보여라.
\begin{align*}
    R(\bm Z) \le \sum_{k=1}^K \log\det\left(\bm I + \alpha\bm Z_k\bm Z_k^T \right),
\end{align*}
여기서 등호는 모든 $k \neq l \in [K]$에 대해 $\bm Z_k^T\bm Z_l = \bm 0$일 때 그리고 그 때만 성립한다.
\medskip 

\noindent 
3. 어떤 $\alpha >0$가 주어졌을 때, 각 $k \in [K]$에 대해 $\alpha_k=m\alpha/m_k$라고 하자. 다음 함수의 1차 임계점에 대한 닫힌 형태를 유도하라:
\begin{align*}
    f(\bm Z_k) = \frac{1}{2}\log\det\left(\bm I + \alpha\bm Z_k\bm Z_k^T \right) - \frac{m_k}{2m}\log\det\left(\bm I + \alpha_k \bm Z_k\bm Z_k^T \right) - \frac{\lambda}{2}\|\bm Z_k\|_F^2. 
\end{align*}
{\em 힌트:} $r_k=\mathrm{rank}(\bm Z_k)$라고 하자. $\bm Z_k$의 다음 특이값 분해를 고려하라:
\begin{align*}
    \bm Z_k = \bm P_k\bm \Sigma_k\bm Q_k^T = \begin{bmatrix}
        \bm P_{k,1} & \bm P_{k,2}
    \end{bmatrix} \begin{bmatrix}
        \tilde{\bm \Sigma}_{k} & \bm 0 \\
        \bm 0 & \bm 0
    \end{bmatrix} 
    \begin{bmatrix}
        \bm Q_{k,1}^T \\ \bm Q_{k,2}^T
    \end{bmatrix}, 
\end{align*}
여기서 $\bm P_k \in \mathcal{O}^d$는 $\bm P_{k,1} \in \R^{d\times r_k}$와 $\bm P_{k,2} \in \R^{d\times (d-r_k)}$를 가지고, $\bm \Sigma_k \in \R^{d\times m_k}$는 $\tilde{\bm \Sigma}_{k} \in \R^{r_k\times r_k}$가 대각 행렬이고, $\bm Q_k \in \mathcal{O}^{m_k}$는 $\bm Q_{k,1} \in \R^{m_k\times r_k}$와 $\bm P_{k,2} \in \R^{m_k \times (m_k-r_k)}$를 갖는다.
\medskip 

\end{exercise}


\begin{exercise}[행렬 역행렬에 대한 노이만 급수]\label{ex:neumannn} 
$\bm A \in \mathbb{R}^{n\times n}$라고 하자. 만약 $\|\bm A\| < 1$이면, 다음을 보여라.
\begin{align}\label{eq:neumann}
    \left( \bm I - \bm A\right)^{-1} = \sum_{k=1}^\infty \bm A^k.  
\end{align}
{\em 힌트:} 증명은 두 단계로 구성된다. \\
(i) {\bf 1단계}: $\|\bm A^k\| \le \|\bm A\|^k$를 사용하여 $\bm A < 1$일 때 무한 급수 $\sum_{k=1}^\infty \bm A^k$가 수렴함을 보여라. \\
(ii) {\bf 2단계}: 행렬 곱 $(\bm I - \bm A) \sum_{k=1}^\infty \bm A^k$를 계산하라.
\end{exercise}


\begin{exercise}
    \eqref{eq1:grad Rcf}와 \eqref{eq2:grad Rcf}의 그래디언트를 계산하라.
\end{exercise}


\begin{exercise}
    $Kp \le d$일 때 \Cref{cor:var_concave_logdet}를 보여라.
\end{exercise}

\end{document}
